\chapter{Heisenberg and lattice vertex algebras}\label{ch:03}
\begingroup\sloppy

\providecommand{\chthreeV}{V}
\providecommand{\chthreeF}{\mathcal{F}}
\providecommand{\chthreeeps}{\varepsilon}
\providecommand{\chthreecoc}{\operatorname{coc}}
\providecommand{\chthreeR}{\mathbb{R}}

This chapter builds, inside Lean, the two families of vertex operator algebras on which the whole
Moonshine construction rests: the \emph{Heisenberg} (free boson) vertex operator algebra of rank
$r$, and the \emph{lattice} vertex operator algebra $V_L$ of an even positive-definite lattice
$L$. Specialising $L$ to the Leech lattice $\Leech$ gives the vertex operator algebra $V_\Leech$,
the starting point of Frenkel--Lepowsky--Meurman's orbifold construction of $\Vnat$. We then
study two kinds of symmetries of lattice vertex operator algebras that the later chapters use
constantly: the \emph{sign automorphisms} $g_\chi$ attached to characters $\chi\colon L\to\Z/2$,
and the \emph{functoriality} of $L\mapsto V_L$ under isometric embeddings of lattices, together with
a rigidity theorem for lifts of the automorphism $-1$ of $L$. The chapter documents the release of the library named in the front matter; every Lean
quotation is checked by \texttt{scripts/textbook/sync.py} to occur verbatim (up to whitespace and
comments) in the source files of that release.

The chapter is written so that a referee can check, for every headline declaration, that the Lean
statement says what the mathematics says. Each definition and theorem is presented in the order:
mathematical statement; verbatim Lean source; a \emph{fidelity check} (in a box headed ``From the
mathematics to the Lean statement''); a proof on paper; and the machine status.

\section{Overview and the physical picture}\label{ch03:sec:overview}

\subsection{What a free boson is}

Physically, a chiral free boson in two dimensions is a field $X(z)$ whose derivative
$h(z)=i\partial X(z)$ has the operator product expansion
\begin{equation}\label{ch03:eq:hhOPE}
  h(z)\,h(w)\sim \frac{k}{(z-w)^2},
\end{equation}
where $k$ is a normalisation constant (the ``level''; one usually takes $k=1$). Expanding
$h(z)=\sum_{n\in\Z} h(n)\,z^{-n-1}$, the OPE \eqref{ch03:eq:hhOPE} is equivalent to the
commutation relations of the \emph{Heisenberg algebra}
\begin{equation}\label{ch03:eq:heis}
  [h(m),h(n)] = m\,k\,\delta_{m+n,0}.
\end{equation}
The state space of the theory is the irreducible representation generated by a vacuum $\mathbb 1$
annihilated by all $h(n)$ with $n\ge 0$: the \emph{Fock space}. Since the negative modes $h(-n)$,
$n\geq1$, commute with one another, the Fock space is simply a polynomial ring in infinitely many
variables, one variable for each $h(-n)$. With $r$ bosons $h_1,\dots,h_r$ (mutually commuting
copies) one gets a polynomial ring in variables indexed by pairs (mode, colour).

If the boson is \emph{compactified}, i.e.\ $X$ takes values in a torus $\mathbb R^r/2\pi L^*$, the
theory acquires \emph{momentum sectors}: states $e^\alpha$ of momentum $\alpha$ in a lattice $L$,
on which the zero modes act by $h(0)e^\alpha=\langle h,\alpha\rangle e^\alpha$. The vertex operator of
a momentum state is the normal-ordered exponential
\begin{equation}\label{ch03:eq:VOphysical}
  Y(e^\alpha,z) = {:}e^{i\alpha\cdot X(z)}{:}
   = \exp\Bigl(\sum_{n\geq1}\frac{\alpha(-n)}{n}z^{n}\Bigr)
     \exp\Bigl(-\sum_{n\geq1}\frac{\alpha(n)}{n}z^{-n}\Bigr)\,e_\alpha\,z^{\alpha(0)} .
\end{equation}
Two such operators satisfy, by a one-line Baker--Campbell--Hausdorff computation that we shall do
in Section~\ref{ch03:sec:locality},
\begin{equation*}
  Y(e^\alpha,z_1)Y(e^\beta,z_2)\;\propto\;(z_1-z_2)^{\langle\alpha,\beta\rangle}\,{:}\cdots{:}\,,
\end{equation*}
so exchanging the two operators produces the sign $(-1)^{\langle\alpha,\beta\rangle}$. If
$\langle\alpha,\beta\rangle$ is odd the two operators \emph{anticommute} rather than commute. To
obtain an honest (bosonic) vertex algebra one inserts ``Klein factors'': the operator $e_\alpha$ is
not simply the shift $e^\beta\mapsto e^{\alpha+\beta}$ but the shift multiplied by a sign
$\varepsilon(\alpha,\beta)$, a $2$-cocycle on $L$ whose antisymmetry
$\varepsilon(\alpha,\beta)/\varepsilon(\beta,\alpha)=(-1)^{\langle\alpha,\beta\rangle}$ exactly
cancels the bad sign. This is the mathematical heart of the lattice construction, and its Lean
version occupies Section~\ref{ch03:sec:cocycle}.

The \emph{evenness} of $L$ (that $\langle\alpha,\alpha\rangle\in 2\Z$) is what makes such a cocycle
exist; it is also what makes all conformal weights $\frac12\langle\alpha,\alpha\rangle$ integral, so
that $V_L$ is $\Z$-graded. \emph{Positive-definiteness} makes every weight space
finite-dimensional, because only finitely many lattice vectors have bounded norm.

\subsection{Road map}

\begin{itemize}[leftmargin=2em]
\item Section~\ref{ch03:sec:conventions} recalls the library's notions of formal
  distribution, mode, field, locality, normal-ordered product, vertex algebra, vertex operator
  algebra, automorphism and homomorphism. These are quoted verbatim and decoded in
  Chapters~\ref{ch:L01}, \ref{ch:L03} and~\ref{ch:L04}; we recall what the statements of this
  chapter use, with pointers.
\item Section~\ref{ch03:sec:heisenberg}: the rank-$r$ Heisenberg vertex operator algebra
  \lean{VOA.HeisenbergN.heisenbergNVOA}.
\item Section~\ref{ch03:sec:cocycle}: integral lattices, the FLM $2$-cocycle
  \lean{latticeCocycle} and its two key properties.
\item Section~\ref{ch03:sec:frame}: how a lattice charge becomes a Heisenberg oscillator, via the
  orthonormal frame $\sqrt G$.
\item Sections~\ref{ch03:sec:latticespace}--\ref{ch03:sec:locality}: the state space
  $V_L=\mathcal F\otimes\C[L]$, its grading, the vertex operators $Y(e^\alpha,z)$, and their mutual
  locality.
\item Section~\ref{ch03:sec:reconstruction}: the reconstruction-from-locality strategy and the
  assembled instance \lean{VOA.LatGen.latticeVOAGen}.
\item Section~\ref{ch03:sec:leech}: the Leech lattice vertex operator algebra
  \lean{VOA.Moonshine.leechVOA}.
\item Section~\ref{ch03:sec:sign}: the sign automorphisms \lean{VOA.LatGen.chiAut}.
\item Section~\ref{ch03:sec:functoriality}: functoriality \lean{VOA.LatGen.latMorph}, the
  $\Z/4$-valued coboundary, character automorphisms and rigidity of lifts of $-1$.
\item Section~\ref{ch03:sec:design}: the design question about the non-computable basis, and why
  it turned out not to matter.
\item Section~\ref{ch03:sec:checklist}: the referee's checklist.
\end{itemize}

\begin{status}
Every headline declaration of this chapter was re-audited for this book with
\texttt{lake env lean -DmaxSynthPendingDepth=3} on a scratch file importing
\leanpath{VOA.Moonshine.LatticeMorphRigid}, \leanpath{VOA.Moonshine.LeechLatticeVOA} and
\leanpath{VOA.Examples.HeisenbergNVOA}. For each of \lean{heisenbergNVOA},
\lean{latticeCocycle_comm}, \lean{latYfullGen_of_locality}, \lean{latticeVOAGen},
\lean{leechVOA}, \lean{leechVOA_centralCharge}, \lean{leechVOA_grading_one}, \lean{chiAut},
\lean{chiLat_eq_self_iff}, \lean{chiAut_mul_thetaAut}, \lean{latticeCocycle_map_mul_latCob},
\lean{latMorph_isVOAHom}, \lean{latMorph_injective}, \lean{mem_range_latMorph_iff},
\lean{latMorph_thetaLat}, \lean{charAut}, \lean{exists_char_of_lift} and
\lean{lift_eq_conj_theta'}, \texttt{\#print axioms} reports exactly
\texttt{[propext, Classical.choice, Quot.sound]}: Lean's three standard axioms. None of them uses
\lean{VOA.pbw} (the PBW theorem of Chapter~\ref{ch:L02}; formerly a project axiom, now a theorem). The only hypotheses are the explicit ones named in each
statement (evenness, positive-definiteness, isometry, equal rank where stated).
\end{status}

\section{Conventions recalled from the foundations}\label{ch03:sec:conventions}

Everything in this section is defined, quoted verbatim from the Lean source, and decoded in the
library chapters; we recall only what the statements of this chapter use, with a pointer to the
place where the exact Lean text and its fidelity check are printed.

\subsection{Formal distributions and modes}

A \emph{formal distribution} \lean{FormalDist 𝒲} (notation \lean{𝒲⟦z^±⟧}) is a doubly infinite
series $\sum_{j\in\Z}f_jz^j$, encoded as the function $j\mapsto f_j$
(Definition~\ref{l01:def:formaldist}, \leanpath{VOA/Foundations/FormalSeries.lean}).
\emph{Warning on indexing}: the argument $j$ of a \lean{FormalDist} is the \emph{exponent of $z$},
not a mode number. Mode numbers are introduced by \lean{mode f m = f (-m-1)}
(Definition~\ref{l01:def:mode}, \leanpath{VOA/Foundations/FormalCalculus.lean}), so that if
$Y(u,z)=\sum_n u_{(n)}z^{-n-1}$ then $u_{(n)}$ is \lean{mode (Y u) n} $=$ \lean{(Y u) (-n-1)}.
Throughout this chapter we shall have to translate between the two indexings; we always say which
one a Lean expression uses. A formal distribution $a$ with values in $\End(V)$ is a \emph{field}
(\lean{IsField a}) if $a(z)x$ has only finitely many negative powers of $z$ for every $x$, i.e.\
$a_{(n)}x=0$ for $n$ large (Definition~\ref{l01:def:field}, \leanpath{VOA/Core/Dong.lean}).

\subsection{Locality and normal ordering}

Two fields are \emph{mutually local}, \lean{Local a b}, if
$(z_1-z_2)^N[a(z_1),b(z_2)]=0$ for some $N\ge0$, where the commutator is the two-variable
distribution \lean{bracketDist a b} (indexed by exponents of $z_1$ and $z_2$) and
\lean{mulBySubPow N} is multiplication by $(z_1-z_2)^N$ (Definitions~\ref{l01:def:formaldist2}
and~\ref{l01:def:bracket}, \leanpath{VOA/Foundations/Locality.lean}). The normal-ordered product
${:}a(z)b(z){:}=a(z)_+b(z)+b(z)a(z)_-$ is the field \lean{nprod ha hb}, built from
\lean{normalOrdered a b}, in which the index $m\ge0$ (the power $z^m$, i.e.\ the creation part
$a(z)_+$) is placed on the left (Definition~\ref{l01:def:field}, \leanpath{VOA/Core/Dong.lean}).

\subsection{Vertex algebras and vertex operator algebras}

The library's \lean{PreVertexAlgebra} (Definition~\ref{l03:def:VA},
\leanpath{VOA/Core/VertexAlgebra.lean}) is a $\Z$-graded space with finite-dimensional pieces
(\lean{grading}, \lean{isInternal}, \lean{finrank}), a linear state--field map
\lean{Y : 𝒱 →ₗ[ℂ] (End ℂ 𝒱)⟦z^±⟧}, a vacuum of degree $0$ with $Y(\mathbb1,z)=\mathrm{id}$, the
creation property, locality, regularity and the grading law $\wt(u_{(n)}v)=\wt u+\wt v-n-1$. The
$n$-th product is \lean{V.smul u n v} $:=$ \lean{mode (V.Y u) n v}. A
\lean{VertexOperatorAlgebra} (Definition~\ref{l03:def:VOA},
\leanpath{VOA/Core/VertexOperatorAlgebra.lean}) adds a \lean{conformalVector} $\omega$ in degree
$2$ and a \lean{centralCharge}, with $L_n=\omega_{(n+1)}=$ \lean{mode (Y conformalVector) (n+1)}
satisfying the Virasoro relations (\lean{virasoro}), $L_0=n$ on \lean{grading n}
(\lean{L0_grading}), $Y(L_{-1}v,z)=\partial_zY(v,z)$ (\lean{translation}), and the CFT-type
conditions $\mathcal V_0=\C\mathbb1$, $\mathcal V_n=0$ for $n<0$ (\lean{vacuum_grade},
\lean{positivity}).

\subsection{Automorphisms, homomorphisms, intertwining}

An automorphism \lean{Aut W} (Definition~\ref{l04:def:aut}, \leanpath{VOA/Core/Automorphism.lean})
is a linear bijection preserving all products \lean{W.smul u n v}, the vacuum and the conformal
vector. A homomorphism \lean{IsVOAHom V T φ} (Definition~\ref{l04:def:hom},
\leanpath{VOA/Core/SimpleImage.lean}) preserves the products and the vacuum but \emph{not
necessarily the conformal vector}. This matters below: an isometric embedding of a lattice of
smaller rank gives a homomorphism that does not send $\omega$ to $\omega$
(Remark~\ref{ch03:rem:notconformal}).

Finally, many proofs in this chapter rest on the notion of an operator $g$ \emph{intertwining} two
formal distributions up to a scalar, $g\,A(z)=c\,B(z)\,g$:
\begin{leancode}
def Intertw (g : End ℂ V) (c : ℂ) (A B : FormalDist (End ℂ V)) : Prop :=
  ∀ (t : ℤ) (x : V), g (A t x) = c • B t (g x)

def IntertwH (g : V →ₗ[ℂ] W) (c : ℂ) (A : FormalDist (End ℂ V)) (B : FormalDist (End ℂ W)) :
    Prop :=
  ∀ (t : ℤ) (x : V), g (A t x) = c • B t (g x)
\end{leancode}
(\leanpath{VOA/Moonshine/LatticeInvolution.lean}, \leanpath{VOA/Moonshine/LatticeMorphFock.lean}),
the second being the version for a map between two different spaces. The basic fact, used
repeatedly, is that intertwining is preserved by normal-ordered products with the scalars
multiplying: if $gA=c\,A'g$ and $gB=d\,B'g$ then $g\,{:}AB{:}=cd\,{:}A'B'{:}\,g$. On paper this is
obvious, because $g$ passes through each factor of each term of the normal-ordered product.

\section{The rank-\texorpdfstring{$r$}{r} Heisenberg vertex operator algebra}\label{ch03:sec:heisenberg}

All declarations of this section live in the namespace \lean{VOA.HeisenbergN} and in the files
\leanpath{VOA/Examples/HeisenbergNFock.lean}, \leanpath{HeisenbergNVertex.lean},
\leanpath{HeisenbergNCentral.lean} and \leanpath{HeisenbergNVOA.lean}.

\subsection{The Fock space as a polynomial ring}

Fix $r\in\N$ and a level $k\in\C$. The rank-$r$ Heisenberg Lie algebra has basis $h_a(n)$
($a=1,\dots,r$, $n\in\Z$) and a central element $K$, with bracket
\begin{equation}\label{ch03:eq:heisr}
  [h_a(m),h_b(n)] = m\,\delta_{ab}\,\delta_{m+n,0}\,K .
\end{equation}
Let $\mathcal F$ be the representation induced from the one-dimensional representation of the
subalgebra spanned by $K$ and the $h_a(n)$, $n\ge 0$, on which $K$ acts by $k$ and the $h_a(n)$ by
$0$. By the Poincar\'e--Birkhoff--Witt theorem (here elementary, since the negative part is
abelian) $\mathcal F$ is the symmetric algebra on the negative modes:
\begin{equation*}
  \mathcal F\;\cong\;\C\bigl[\,h_a(-n)\;:\;1\le a\le r,\;n\ge 1\,\bigr]\cdot \mathbb 1 .
\end{equation*}
The library \emph{defines} the Fock space to be this polynomial ring, rather than constructing it
as an induced module. Variables are indexed by pairs $(i,a)\in\N\times\{0,\dots,r-1\}$, and the
variable $X_{(i,a)}$ stands for the state $h_a(-i-1)\mathbb 1$. As $i$ runs through $0,1,2,\ldots$ this gives all
creation modes $h_a(-1),h_a(-2),h_a(-3),\ldots$; indexing by $i=-n-1\in\N$ rather than by the mode
$n\le-1$ avoids negative integers in the variable names.

\begin{definition}[Fock space]\label{ch03:def:FockN}
$\mathcal F_r := \C[X_{(i,a)} : i\in\N,\ a\in\{0,\dots,r-1\}]$, the polynomial ring in countably
many variables. The vacuum is the constant polynomial $1$.
\end{definition}

\begin{leancode}
abbrev FockN (r : ℕ) : Type := MvPolynomial (ℕ × Fin r) ℂ

noncomputable def vacuumN (r : ℕ) : FockN r := (1 : FockN r)
\end{leancode}

\begin{dictionary}
\lean{MvPolynomial σ ℂ} is Mathlib's ring of polynomials with complex coefficients in variables
indexed by the type $\sigma$; here $\sigma=\N\times\lean{Fin r}$, where \lean{Fin r} is the
$r$-element type $\{0,\dots,r-1\}$ (colours are numbered from $0$). The variable for $(i,a)$ is
written \lean{X (i, a)}. So \lean{X (0, a)} is $h_a(-1)\mathbb 1$, \lean{X (1, a)} is
$h_a(-2)\mathbb 1$, and the monomial \lean{X (0,a)^2 * X (3,b)} is
$h_a(-1)^2h_b(-4)\mathbb 1$. The identification of $\mathcal F$ with the induced module is a
theorem of mathematics, not of the library; what the library needs, and proves, is that the
oscillators defined below satisfy the Heisenberg relations on this space and that the vacuum is
annihilated by the non-negative modes. Those two facts characterise the induced module up to
isomorphism (it is irreducible for $k\ne 0$), so nothing is lost. This definition cannot be
vacuous: a polynomial ring is a perfectly concrete, nonzero, infinite-dimensional space.
\end{dictionary}

\begin{physics}
The Fock space is the Hilbert space (algebraically: the space of finite-energy states) of $r$ free chiral bosons $\partial X^a$: polynomials in the variables $X_{(i,a)}$ are the states obtained from the vacuum by applying creation operators, with $X_{(i,a)}$ standing for the oscillator $\alpha^a_{-(i+1)}$. The vacuum $1$ is the $SL_2$-invariant state with no oscillator excited, and the neutral sector (zero momentum) is exactly what this space describes. Heuristically it is the oscillator part of the closed-string Hilbert space on a circle with the zero modes removed; the definition itself is just a polynomial ring.
\end{physics}

\subsection{The oscillators}

The creation operator $h_a(-m)$, $m\ge1$, is multiplication by the variable $X_{(m-1,a)}$. The
annihilation operator $h_a(n)$, $n\ge 1$, must satisfy $[h_a(n),h_a(-n)]=nk$; since
$[\partial/\partial X,X]=1$, the right choice is $h_a(n)=nk\,\partial/\partial X_{(n-1,a)}$. The
zero mode acts by $0$ on the Fock space (it is the momentum, which vanishes in the neutral sector).

\begin{definition}[Oscillators]\label{ch03:def:oscN}
For $a\in\{0,\dots,r-1\}$ and $n\in\Z$ define $h_a(n)\in\End\mathcal F_r$ by
\begin{equation*}
  h_a(n)=\begin{cases} X_{(-n-1,a)}\cdot(-) & n<0,\\ 0 & n=0,\\
  n\,k\,\dfrac{\partial}{\partial X_{(n-1,a)}} & n>0.\end{cases}
\end{equation*}
\end{definition}

\noindent File \leanpath{VOA/Examples/HeisenbergNFock.lean}.
\begin{leancode}
noncomputable def mulXN {r : ℕ} (i : ℕ) (a : Fin r) : FockN r →ₗ[ℂ] FockN r :=
  (LinearMap.mul ℂ (FockN r) (X (i, a))).comp (LinearMap.id)

noncomputable def pderivLMN {r : ℕ} (i : ℕ) (a : Fin r) : FockN r →ₗ[ℂ] FockN r :=
  (pderiv (i, a) : Derivation ℂ (FockN r) (FockN r)).toLinearMap

noncomputable def oscN (r : ℕ) (k : ℂ) (a : Fin r) (n : ℤ) : FockN r →ₗ[ℂ] FockN r :=
  if n < 0 then mulXN (-n - 1).toNat a
  else if 0 < n then (n : ℂ) • k • pderivLMN (n - 1).toNat a
  else 0
\end{leancode}

\begin{dictionary}
\lean{→ₗ[ℂ]} denotes $\C$-linear maps. \lean{pderiv (i,a)} is Mathlib's partial derivative with
respect to the variable \lean{X (i,a)}, a derivation of the polynomial ring. \lean{(-n-1).toNat}
converts the integer $-n-1$ to a natural number (truncating negative values to $0$); in the branch
$n<0$ the integer $-n-1$ is $\ge 0$, so no truncation happens and the variable is $X_{(-n-1,a)}$,
exactly as in Definition~\ref{ch03:def:oscN}. Similarly \lean{(n-1).toNat} $=n-1$ in the branch
$n>0$. The scalar is $n\cdot k$, written \lean{(n : ℂ) • k • ...}. The Lean definition is
literally the textbook one.
\end{dictionary}

\begin{physics}
The operators $h_a(n)$ are the Laurent modes $\alpha^a_n$ of the chiral current $i\partial X^a(z)=\sum_n\alpha^a_nz^{-n-1}$ of a free boson. For $n<0$ they create a quantum of energy $|n|$ (multiplication by a variable), for $n>0$ they annihilate one (a derivative), and the zero mode is the momentum, which vanishes on the neutral sector and is reinstated for lattice states later. The normalisation $nk$ in the annihilators is chosen so that the algebra is the standard oscillator algebra of the next theorem, with $k$ the level (normalisation of the two-point function $\langle J(z)J(w)\rangle=k/(z-w)^2$).
\end{physics}

\begin{theorem}[Heisenberg relations]\label{ch03:thm:oscN_comm}
For all colours $a,b$ and all $m,n\in\Z$,
$[h_a(m),h_b(n)] = m\,k\,\delta_{ab}\,\delta_{m+n,0}\cdot\mathrm{id}$ on $\mathcal F_r$.
\end{theorem}

\noindent File \leanpath{VOA/Examples/HeisenbergNFock.lean}.
\begin{leancode}
theorem oscN_comm (r : ℕ) (k : ℂ) (a b : Fin r) (m n : ℤ) :
    oscN r k a m * oscN r k b n - oscN r k b n * oscN r k a m
      = (if a = b ∧ m + n = 0 then (m : ℂ) * k else 0) • LinearMap.id
\end{leancode}

\begin{dictionary}
The product \lean{*} of two endomorphisms is composition. The right-hand side is the scalar
$mk$ if $a=b$ and $m+n=0$, and $0$ otherwise, times the identity map: this is
$mk\,\delta_{ab}\delta_{m+n,0}$. No hypothesis is imposed on $k$ (it may even be $0$, in which case
all annihilators vanish and the relation still holds). The statement is the textbook one,
quantified over all integers $m,n$, including the zero modes.
\end{dictionary}

\begin{physics}
This is the canonical commutation relation of the free chiral boson, $[\alpha_m,\alpha_n]=m\delta_{m+n,0}$ (for $k=1$): an infinite family of harmonic oscillators, one for each positive mode number. It is the algebraic form of the current--current OPE $J^a(z)J^b(w)\sim k\,\delta^{ab}/(z-w)^2$, whose double pole is the entire singular part. Since all the physics of the free boson is encoded in these relations, the Fock space is the unique irreducible module in which the $h_a(-m)$ act freely.
\end{physics}

\begin{proof}[Proof on paper]
Split according to the signs of $m$ and $n$. Two creation operators are multiplications by
variables, which commute; and $m+n\ne 0$, so the right side is $0$. Two annihilation operators
are multiples of partial derivatives, which commute; again $m+n\neq0$. A zero mode is $0$ and so
commutes with everything, and $m\cdot\delta_{m+n,0}$ vanishes when $m=0$, while if $n=0\ne m$ the
delta vanishes. Finally take $m<0<n$: then $h_a(m)$ is multiplication by $X=X_{(-m-1,a)}$ and
$h_b(n)$ is $nk\,\partial/\partial Y$ with $Y=X_{(n-1,b)}$. If $X\neq Y$ the two commute
(different variables), and $X\ne Y$ happens exactly when $a\ne b$ or $m+n\ne0$. If $X=Y$ the
Leibniz rule gives $\partial_X(Xp)-X\partial_Xp=p$, so
$[X\cdot, nk\,\partial_X]=-nk=mk$ since $n=-m$. The case $m>0>n$ is symmetric.
\end{proof}

\begin{status}
Proved, unconditional, in \leanpath{VOA/Examples/HeisenbergNFock.lean}; the file documents
\texttt{\#print axioms} $=$ \texttt{[propext, Classical.choice, Quot.sound]}.
\end{status}

\begin{proposition}[Vacuum]\label{ch03:prop:vacuum}
$h_a(n)\mathbb 1=0$ for every colour $a$ and every $n\geq 0$.
\end{proposition}
\begin{leancode}
theorem oscN_vacuum (r : ℕ) (k : ℂ) (a : Fin r) (n : ℤ) (hn : 0 ≤ n) :
    oscN r k a n (vacuumN r) = 0
\end{leancode}

\begin{physics}
This says the Fock vacuum is the ground state: it is annihilated by every annihilation operator and by the momentum zero mode. In field language it is the statement that the current is regular at the origin acting on the vacuum, $J(z)\mathbb 1=\sum_{n\le-1}h(n)\mathbb 1\,z^{-n-1}$, i.e.\ the state--operator correspondence sends $\mathbb 1$ to the identity operator. Technically it is the oscillator vacuum condition used in every normal-ordering argument.
\end{physics}
\begin{proof}[Proof on paper]
The derivative of a constant is $0$, and the zero mode is $0$.
\end{proof}

\begin{example}[A small computation]\label{ch03:ex:smallcomp}
Take $r=1$, $k=1$ and write $x_i=X_{(i,0)}$. The state $h(-1)^2\mathbb 1$ is $x_0^2$, and
$h(1)x_0^2 = 1\cdot\partial_{x_0}(x_0^2)=2x_0 = 2h(-1)\mathbb 1$, in agreement with the textbook
computation $h(1)h(-1)^2\mathbb 1=[h(1),h(-1)]h(-1)\mathbb1+h(-1)[h(1),h(-1)]\mathbb1=2h(-1)\mathbb1$.
Likewise $h(2)x_1=2\,\partial_{x_1}x_1=2$, i.e.\ $h(2)h(-2)\mathbb 1=2\mathbb 1$.
\end{example}

\subsection{Currents}

The basic current of colour $a$ is $h_a(z)=\sum_n h_a(n)z^{-n-1}$. As a \lean{FormalDist}
(indexed by exponents of $z$) its $z^j$-coefficient is $h_a(-j-1)$:
\begin{leancode}
noncomputable def currentN (r : ℕ) (k : ℂ) (a : Fin r) :
    FormalDist (Module.End ℂ (FockN r)) := fun j => oscN r k a (-j - 1)

theorem currentN_mode (r : ℕ) (k : ℂ) (a : Fin r) (m : ℤ) :
    mode (currentN r k a) m = oscN r k a m

theorem currentN_local (r : ℕ) (k : ℂ) (a b : Fin r) :
    Local (currentN r k a) (currentN r k b)
\end{leancode}
\lean{currentN_mode} confirms the dictionary: the $m$-th mode (in the sense of \lean{mode}) of
the current is the oscillator $h_a(m)$. \lean{currentN_local} is the locality
$(z_1-z_2)^2[h_a(z_1),h_b(z_2)]=0$; it is proved by computing the commutator as
$k\delta_{ab}\,\partial_{z_2}\delta(z_1/z_2)$ (Theorem~\ref{ch03:thm:oscN_comm} in generating-function
form) and using that $(z_1-z_2)^2$ kills the derivative of the formal delta function.

A trap worth recording: in the definition of the currents the \emph{Lean variable} \lean{j} is a
power of $z$, while in \lean{oscN} the integer \lean{n} is a mode. Mixing up the two produces
off-by-one errors in every formula; the library's lemma \lean{currentN_mode} pins the convention
down once and for all.

\subsection{The conformal grading}

The conformal weight of $h_a(-i-1)\mathbb 1$ is $i+1$, and weights add under multiplication.

\begin{definition}[Weight spaces]\label{ch03:def:gradingWN}
Give the variable $X_{(i,a)}$ weight $i+1$. For $w\in\N$ let $(\mathcal F_r)_w$ be the span of
monomials of total weight $w$.
\end{definition}
\noindent File \leanpath{VOA/Examples/HeisenbergNVertex.lean}.
\begin{leancode}
def weightFnN {r : ℕ} : ℕ × Fin r → ℕ := fun u => u.1 + 1

def gradingWN (r : ℕ) (w : ℕ) : Submodule ℂ (FockN r) :=
  weightedHomogeneousSubmodule ℂ weightFnN w

noncomputable def gradingZN (r : ℕ) (n : ℤ) : Submodule ℂ (FockN r) :=
  if 0 ≤ n then gradingWN r n.toNat else ⊥
\end{leancode}

\begin{dictionary}
\lean{weightedHomogeneousSubmodule ℂ w d} is Mathlib's submodule of polynomials all of whose
monomials have $w$-weighted degree $d$; with $w(i,a)=i+1$ (the field \lean{u.1} is the first
component $i$) this is exactly $(\mathcal F_r)_d$. The $\Z$-indexed version \lean{gradingZN} is
needed because a vertex algebra's grading is indexed by $\Z$: it is $(\mathcal F_r)_n$ for $n\ge0$
and the zero submodule \lean{⊥} for $n<0$.
\end{dictionary}

\begin{physics}
The weight of a state is its $L_0$ eigenvalue, the energy on the cylinder measured from the Casimir shift, and here it counts the total oscillator level: $X_{(i,a)}$ carries energy $i+1$. A weight-$w$ state is a multi-particle state of $r$ species of bosons with total energy $w$, so $\dim(\mathcal F_r)_w$ is the number of $r$-coloured partitions of $w$. Summing gives the oscillator partition function $\operatorname{Tr}q^{L_0}=\prod_n(1-q^n)^{-r}$, which is $q^{r/24}/\eta(\tau)^{r}$ once the Casimir energy is restored. That the weight-$0$ space is spanned by the vacuum and every weight space is finite is the finiteness of the spectrum (a gapped, discrete spectrum for each energy), a basic requirement of a unitary chiral CFT.
\end{physics}

The weight-$w$ space has a basis of monomials $\prod X_{(i,a)}^{e_{(i,a)}}$ with
$\sum e_{(i,a)}(i+1)=w$, i.e.\ of $r$-coloured partitions of $w$. Hence
\begin{equation}\label{ch03:eq:fockchar}
  \sum_{w\ge0}\dim(\mathcal F_r)_w\,q^w=\prod_{n\ge1}(1-q^n)^{-r}.
\end{equation}
Indeed, a monomial is a choice of an independent exponent $e\ge0$ for each variable $X_{(i,a)}$,
contributing weight $e(i+1)$, so the generating function factorises as
$\prod_{(i,a)}\sum_{e\ge0}q^{e(i+1)}=\prod_{a}\prod_{n\ge1}(1-q^n)^{-1}$, one geometric series per
variable.
For $r=1$ the dimensions are $1,1,2,3,5,7,\dots$ (the partition numbers). For $r=24$ the first are
$1$, $24$, $24+\binom{25}{2}=324$, $3200$. Finite-dimensionality of each weight space is proved in
the library (\lean{gradingWN_finiteDimensional}, \lean{gradingZN_finrank}), as is the fact that
$(\mathcal F_r)_0=\C\cdot 1$ (\lean{gradingWN_zero_eq_span_one}). Creation modes $h_a(-m)$ raise
weight by $m$ and annihilation modes $h_a(m)$ lower it by $m$ (\lean{oscN_neg_weight},
\lean{oscN_pos_weight}).

\subsection{Vertex operators and the vertex algebra}

The vertex operator of a monomial state is the normal-ordered product of derivatives of the
currents:
\begin{equation}\label{ch03:eq:YFock}
  Y\bigl(h_{a_1}(-i_1-1)\cdots h_{a_s}(-i_s-1)\mathbb 1,z\bigr)
   = {:}\,\partial^{(i_1)}h_{a_1}(z)\cdots\partial^{(i_s)}h_{a_s}(z)\,{:},
  \qquad \partial^{(i)}:=\tfrac{1}{i!}\partial_z^i .
\end{equation}
The library implements \eqref{ch03:eq:YFock} as follows: \lean{genFieldN r k a i} is
$\frac1{i!}\partial^ih_a(z)$; \lean{YwordN r k l} is the iterated normal-ordered product
${:}g_1{:}g_2{:}\cdots{:}g_s\,\mathrm{id}{:}\cdots{:}$ over a list $l$ of indices (nested to the
right); \lean{monoWordN} turns an exponent vector into such a list; and $Y$ is extended linearly:
\noindent File \leanpath{VOA/Examples/HeisenbergNVertex.lean}.
\begin{leancode}
noncomputable def genFieldN (r : ℕ) (k : ℂ) (a : Fin r) (i : ℕ) :
    FormalDist (End ℂ (FockN r)) :=
  (i.factorial : ℂ)⁻¹ • dcurrentN r k a i

noncomputable def monoWordN (a : (ℕ × Fin r) →₀ ℕ) : List (ℕ × Fin r) :=
  a.support.toList.flatMap (fun ia => List.replicate (a ia) ia)

noncomputable def YmonN (r : ℕ) (k : ℂ) (a : (ℕ × Fin r) →₀ ℕ) :
    FormalDist (End ℂ (FockN r)) := YwordN r k (monoWordN a)

noncomputable def YN (r : ℕ) (k : ℂ) : FockN r →ₗ[ℂ] FormalDist (End ℂ (FockN r)) :=
  Finsupp.linearCombination ℂ (YmonN r k)
\end{leancode}

\begin{dictionary}
\lean{dcurrentN r k a j} is the $j$-th formal derivative $\partial_z^jh_a(z)$. An exponent vector
\lean{(ℕ × Fin r) →₀ ℕ} is a finitely supported function assigning to each variable its exponent;
Mathlib represents the monomial with exponent vector $e$ as \lean{monomial e 1}, and a polynomial
is a finitely supported function from exponent vectors to coefficients, so
\lean{Finsupp.linearCombination ℂ (YmonN r k)} is the linear map sending
$\sum_e c_e\,x^e$ to $\sum_e c_e\,Y_{\mathrm{mon}}(e)$. The list \lean{monoWordN} enumerates the
support in an unspecified order (\lean{Finset.toList}); this is harmless because the fields
$\partial^{(i)}h_a$ are mutually local with commutators that are scalars, so the iterated
normal-ordered products of the free-boson currents do not depend on the order of the factors. A
referee may worry that such an arbitrary choice could make $Y$ ``wrong'' on some monomial; it
cannot, because the creation property proved below (\lean{YN_creation}) pins $Y(u,z)\mathbb 1$ to
$u$, and Goddard's uniqueness theorem then pins $Y(u,z)$ itself among local fields.
\end{dictionary}

The five vertex algebra axioms are then proved one by one (\lean{YN_vacuum}, \lean{YN_creation},
\lean{YN_local}, \lean{va5N_regularity}, \lean{va1N_grading_mode}) and packaged:
\noindent File \leanpath{VOA/Examples/HeisenbergNVOA.lean}.
\begin{leancode}
noncomputable def heisVAN (r : ℕ) (k : ℂ) (_hk : k ≠ 0) : PreVertexAlgebra (FockN r) where
  grading := gradingZN r
  isInternal := gradingZN_isInternal r
  finrank := gradingZN_finrank r
  Y := YN r k
  vacuum := vacuumN r
  ...
\end{leancode}
(the elided fields are the axiom proofs just named). Locality of arbitrary monomial fields is
obtained from locality of the generating currents by \emph{Dong's lemma}: if $a,b,c$ are pairwise
local, then ${:}ab{:}$ is local with $c$ (Theorem~\ref{l03:thm:dong}).

\subsection{The Sugawara conformal vector and the Virasoro relation}

For $k\neq0$, the conformal vector of the free boson is
\begin{equation}\label{ch03:eq:omegaHeis}
  \omega=\frac{1}{2k}\sum_{a}h_a(-1)^2\mathbb 1,\qquad
  L_m=\frac1{2k}\sum_a\sum_{j\in\Z}{:}h_a(m-j)h_a(j){:} .
\end{equation}
\noindent File \leanpath{VOA/Examples/HeisenbergNCentral.lean}.
\begin{leancode}
noncomputable def ωNa (r : ℕ) (k : ℂ) (a : Fin r) : FockN r :=
  (1 / (2 * k)) • oscN r k a (-1) (oscN r k a (-1) (vacuumN r))

noncomputable def ωN (r : ℕ) (k : ℂ) : FockN r := ∑ a : Fin r, ωNa r k a

noncomputable def LvirNa (r : ℕ) (k : ℂ) (a : Fin r) (m : ℤ) : End ℂ (FockN r) :=
  (1 / (2 * k)) • nprod (currentN_isField r k a) (currentN_isField r k a) (-m - 2)

noncomputable def LvirN (r : ℕ) (k : ℂ) (m : ℤ) : End ℂ (FockN r) :=
  ∑ a : Fin r, LvirNa r k a m
\end{leancode}

\begin{dictionary}
The field ${:}h_a(z)h_a(z){:}$ is \lean{nprod ... ...}; its coefficient of $z^{-m-2}$ is the
mode $\sum_j{:}h_a(m-j)h_a(j){:}$ of a weight-$2$ field, i.e.\ $2k\,L^a_m$. So \lean{LvirNa r k a m}
is $L^a_m$ as in \eqref{ch03:eq:omegaHeis}. The library later proves
$L_m=\omega_{(m+1)}$ (\lean{LvirN_eq_mode_YN_omega}), connecting these explicit operators to the
axioms of a vertex operator algebra, which are phrased in terms of $Y(\omega,z)$.
\end{dictionary}

\begin{theorem}[Virasoro relation, $c=r$]\label{ch03:thm:virasoroN}
For $k\ne0$ and all $m,n\in\Z$,
\begin{equation*}
[L_m,L_n]=(m-n)L_{m+n}+\delta_{m+n,0}\,\frac{r(m^3-m)}{12}\,\mathrm{id}.
\end{equation*}
\end{theorem}
\begin{leancode}
theorem L_virasoroN (r : ℕ) (k : ℂ) (hk : k ≠ 0) (m n : ℤ) :
    LvirN r k m * LvirN r k n - LvirN r k n * LvirN r k m
      = ((m : ℂ) - n) • LvirN r k (m + n)
        + (if m + n = 0 then (r : ℂ) * ((m : ℂ)^3 - m) / 12 else 0)
            • (LinearMap.id : End ℂ (FockN r))
\end{leancode}

\begin{physics}
This is the statement that the Sugawara stress tensor $T=\frac1{2k}\sum_a{:}J^aJ^a{:}$ of $r$ free bosons generates a Virasoro algebra with central charge $c=r$: each free boson contributes $c=1$. It is the classical Sugawara construction in the abelian case and is verified by the same Wick-contraction computation that physicists use to derive the central term $\frac{c}{12}(m^3-m)$ from the double contraction of two currents. The result is independent of the level $k$, since $T$ is quadratic in $J$ and rescales accordingly.
\end{physics}

\begin{proof}[Proof on paper]
First, $[L^a_m,h_b(n)]=-n\,\delta_{ab}\,h_b(m+n)$: commute $h_b(n)$ through each normal-ordered
quadratic term using \eqref{ch03:eq:heisr}. For $a\neq b$ the operators $L^a_m$ and $L^b_n$ commute,
because they are built from mutually commuting oscillators. For $a=b$ one computes
$[L^a_m,L^a_n]$ by applying $[L^a_m,\cdot\,]$ to the defining sum of $L^a_n$; the result is
$(m-n)L^a_{m+n}$ plus a term that arises from re-normal-ordering, which is a finite ``window'' sum
$\frac{1}{2}\sum_{j}(\ldots)$ that evaluates to $\frac{m^3-m}{12}\delta_{m+n,0}$, the familiar $c=1$
anomaly of one free boson. (The library computes this window by the explicit power sums
$\sum i=N(N-1)/2$ and $\sum i^2=N(N-1)(2N-1)/6$.) Summing over the $r$ colours gives $c=r$.
\end{proof}

\begin{status}
Proved, conditional only on $k\neq0$ (needed to divide by $2k$), in
\leanpath{VOA/Examples/HeisenbergNVOA.lean}.
\end{status}

\subsection{The instance}

\begin{theorem}[Heisenberg vertex operator algebra]\label{ch03:thm:heisVOA}
For every $r\in\N$ and every $k\in\C\setminus\{0\}$, the Fock space $\mathcal F_r$ with the
vertex operators \eqref{ch03:eq:YFock}, vacuum $1$, conformal vector \eqref{ch03:eq:omegaHeis}
and grading by weight is a vertex operator algebra of CFT type with central charge $r$.
\end{theorem}
\begin{leancode}
noncomputable def heisenbergNVOA (r : ℕ) (k : ℂ) (hk : k ≠ 0) :
    VertexOperatorAlgebra (FockN r) where
  toPreVertexAlgebra := heisVAN r k hk
  conformalVector := ωN r k
  conformal_mem := ωN_mem_gradingZN_two r k
  centralCharge := (r : ℂ)
  virasoro := fun m n => conf_virasoroN r k hk m n
  L0_grading := fun v hv => conf_L0_gradingN r k hk v hv
  translation := fun v => conf_translationN r k hk v
  vacuum_grade := gradingZN_zero_eq_span_vacuum r
  positivity := fun n hn => gradingZN_positivity r n hn

theorem heisenbergN_central_charge (r : ℕ) (k : ℂ) (hk : k ≠ 0) :
    (heisenbergNVOA r k hk).centralCharge = (r : ℂ) := rfl
\end{leancode}

\begin{dictionary}
A definition of a structure value is a theorem in disguise: to write down
\lean{heisenbergNVOA} Lean demands a proof of every axiom field, and these are the named lemmas.
The only hypothesis is \lean{hk : k ≠ 0}. The central charge lemma is proved by \lean{rfl} (it
just reads off the field \lean{centralCharge}); this is \emph{not} a vacuous statement, because
the field \lean{virasoro} forces the central charge to be what it is. Indeed $[L_2,L_{-2}]=4L_0+
\frac{c}{2}\mathrm{id}$ applied to the vacuum gives $L_2\omega=\frac c2\mathbb 1$, and
$L_2\omega=\frac r2\mathbb 1$ by direct computation; any wrong value of \lean{centralCharge} would
make the \lean{virasoro} field unprovable. \lean{translation} is proved via Goddard's uniqueness
theorem and the covariance $[L_{-1},Y(v,z)]=\partial_zY(v,z)$ (\lean{YN_DCovariant}).
\end{dictionary}

\begin{physics}
This theorem says that the free boson theory is a consistent chiral conformal field theory: the Fock space is the space of states, the normal-ordered products of derivatives of the current are the local fields, the vertex operator of $h_{-1}\mathbb 1$ is the current $J(z)$, and the conformal vector gives the stress tensor with central charge $c=r$. Locality of the fields is the physical statement that correlators are single-valued up to meromorphic singularities, and CFT type that the spectrum is bounded below with a unique vacuum. The associated partition function is $q^{-r/24}\prod(1-q^n)^{-r}=1/\eta^r$ up to the factor $q^{r/24}$. This is the model that the Moonshine construction then enlarges by adding momentum.
\end{physics}

\begin{status}
Proved; \texttt{\#print axioms VOA.HeisenbergN.heisenbergNVOA} $=$
\texttt{[propext, Classical.choice, Quot.sound]} (re-checked for this book). Conditional only on
$k\neq 0$. The lattice vertex operator algebra uses the level $k=1$.
\end{status}

\section{Even lattices and the FLM 2-cocycle}\label{ch03:sec:cocycle}

\subsection{Integral lattices in the library}

Lattices are treated in detail in Chapter~\ref{ch:02}. The vertex-algebra side of the library uses
the following minimal notion.

\begin{definition}[Integral lattice]\label{ch03:def:lattice}
An \emph{integral lattice} is a finitely generated free abelian group $L$ together with a
symmetric $\Z$-bilinear form $\langle\cdot,\cdot\rangle\colon L\times L\to\Z$. It is \emph{even}
if $\langle v,v\rangle\in2\Z$ for all $v$, and \emph{positive-definite} if
$\langle v,v\rangle>0$ for all $v\ne0$. Its \emph{rank} is the rank of $L$ as a free
$\Z$-module.
\end{definition}

The Lean structure \lean{VOA.IntegralLattice} and the predicates \lean{Even}, \lean{PosDef} and
\lean{rank} are quoted verbatim and decoded in Definitions~\ref{ch02:def:voalattice},
\ref{ch02:def:voaevenposdef}, \ref{ch02:def:voarank} and~\ref{l02:def:lattice}
(\leanpath{VOA/Foundations/Lattices.lean}, \leanpath{VOA/Foundations/LatticeCocycle.lean}). We
recall that the carrier is a free, finitely generated $\Z$-module but \emph{no basis is part of the
data}, and that \lean{Even} and \lean{PosDef} appear explicitly, as \lean{hEven} and \lean{hPos},
in every theorem that needs them.

\begin{dictionary}
Positivity is asked only on integer vectors; that it implies positivity of the real extension
(and therefore finiteness of $\{v:\langle v,v\rangle\le N\}$) is a theorem of the library,
\lean{IntegralLattice.boundedFinite} in \leanpath{VOA/Foundations/LatticeFinite.lean}:
\begin{leancode}
def BoundedFinite (L : IntegralLattice) : Prop :=
  ∀ N : ℤ, {α : L.carrier | L.form α α ≤ N}.Finite

theorem boundedFinite (hPos : L.PosDef) : BoundedFinite L
\end{leancode}
Note that \lean{PosDef} on integer vectors is genuinely equivalent to real positive-definiteness:
a rational positive quadratic form is positive on $\Q^n$, hence non-negative on $\chthreeR^n$ by
density, and a non-degenerate non-negative form is positive; non-degeneracy over $\Q$ follows from
positivity on $\Q^n$. So nothing is lost.
\end{dictionary}

\begin{physics}
An even positive-definite lattice $L$ of rank $r$ is the data of a compactification: the target space is the torus $\mathbb R^r/2\pi L^*$, and the vectors of $L$ are the allowed (chiral) momenta. Evenness, $\langle\alpha,\alpha\rangle\in2\Z$, says that the vertex operator $:\!e^{i\alpha\cdot X}\!:$ has integer conformal weight $\langle\alpha,\alpha\rangle/2$, so that it is a bosonic single-valued field; integrality of the pairing makes two such operators mutually local. Positive-definiteness is what makes the momentum $p\mapsto p^2/2$ energy bounded below. Heuristically this is the standard dictionary lattice $\leftrightarrow$ chiral algebra of a compact torus at a given background; the definition is the algebraic datum alone.
\end{physics}

\subsection{Why a cocycle is needed}

Let $L$ be even. As explained in Section~\ref{ch03:sec:overview}, the naive vertex operators
$Y(e^\alpha,z)$ built with the untwisted shift $e^\beta\mapsto e^{\alpha+\beta}$ satisfy
$Y(e^\alpha,z_1)Y(e^\beta,z_2)=(-1)^{\langle\alpha,\beta\rangle}Y(e^\beta,z_2)Y(e^\alpha,z_1)$
``up to the $(z_1-z_2)^N$ ambiguity''. The cure (Frenkel--Kac, Segal; FLM Chapter~5 and~8) is to
replace the group algebra $\C[L]$ by a twisted group algebra $\C_\varepsilon[L]$, with product
$e^\alpha\cdot e^\beta=\varepsilon(\alpha,\beta)e^{\alpha+\beta}$, where
$\varepsilon\colon L\times L\to\{\pm1\}$ satisfies
\begin{align}
  \varepsilon(\alpha,\beta)\varepsilon(\alpha+\beta,\gamma)&=\varepsilon(\beta,\gamma)
    \varepsilon(\alpha,\beta+\gamma) &&\text{(2-cocycle law: associativity)},\label{ch03:eq:coclaw}\\
  \varepsilon(\alpha,\beta)\varepsilon(\beta,\alpha)&=(-1)^{\langle\alpha,\beta\rangle}
  &&\text{(commutator property)}.\label{ch03:eq:commprop}
\end{align}
A convenient way to produce such an $\varepsilon$ is to make it \emph{bimultiplicative}:
$\varepsilon(\alpha,\beta)=(-1)^{B(\alpha,\beta)}$ for a $\Z$-bilinear form $B$. Then
\eqref{ch03:eq:coclaw} is automatic, and \eqref{ch03:eq:commprop} asks that
$B(\alpha,\beta)+B(\beta,\alpha)\equiv\langle\alpha,\beta\rangle\pmod 2$. Choosing a basis
$B_1,\dots,B_n$ of $L$ one takes the strictly triangular part of the Gram matrix.

\subsection{The cocycle in the library}

\begin{definition}[Cocycle exponent and cocycle]\label{ch03:def:cocExp}
Let $\mathcal B=(B_i)_{i<n}$ be a $\Z$-basis of $L$ and write $\alpha=\sum_i\alpha_iB_i$. Put
\begin{equation*}
  \chthreecoc_{\mathcal B}(\alpha,\beta)=\sum_{j<i}\langle B_i,B_j\rangle\,\alpha_i\beta_j,\qquad
  \varepsilon_{\mathcal B}(\alpha,\beta)=(-1)^{\chthreecoc_{\mathcal B}(\alpha,\beta)}\in\C .
\end{equation*}
\end{definition}
\begin{leancode}
noncomputable def cocExp (α β : L.carrier) : ℤ :=
  ∑ i : Fin n, ∑ j : Fin n,
    (if j < i then L.form (B i) (B j) else 0) * B.repr α i * B.repr β j

noncomputable def cocycle (α β : L.carrier) : ℂ := (-1 : ℂ) ^ (cocExp L B α β)
\end{leancode}
(\leanpath{VOA/Foundations/LatticeCocycle.lean}, with the section variable
\lean{(B : Basis (Fin n) ℤ L.carrier)}.)

\begin{dictionary}
\lean{B.repr α i} is the $i$-th coordinate $\alpha_i\in\Z$ of $\alpha$ in the basis $B$. The sum
over \lean{i j : Fin n} with the guard \lean{j < i} is the sum over pairs with $j<i$. The power
\lean{(-1 : ℂ) ^ (k : ℤ)} is an integer power of $-1$ in $\C$; the cocycle takes values $\pm1$
viewed as complex numbers, because it will multiply complex vectors.
\end{dictionary}

\begin{physics}
The sign $\varepsilon(\alpha,\beta)$ is the Klein-factor (cocycle) correction in the product of vertex operators $:\!e^{i\alpha\cdot X}\!:$. Naively these exponentials commute up to the phase $(-1)^{\langle\alpha,\beta\rangle}$ coming from the OPE $(z-w)^{\langle\alpha,\beta\rangle}$, and for odd $\langle\alpha,\beta\rangle$ they would anticommute where bosons should commute. Multiplying the operators by momentum-shifting signs $\varepsilon$ compensates, exactly as Klein factors in bosonisation convert statistics. The particular choice via the triangular part of the Gram matrix is a gauge: physical quantities are independent of it up to isomorphism, and the cocycle is the cohomology class in $H^2(L,\Z/2)$ fixed by the form.
\end{physics}

\begin{theorem}[Properties of $\varepsilon_{\mathcal B}$]\label{ch03:thm:cocycle}
For every basis $\mathcal B$ and all $\alpha,\alpha',\beta,\beta',\gamma\in L$:
\begin{enumerate}[label=(\roman*)]
\item $\varepsilon(\alpha+\alpha',\beta)=\varepsilon(\alpha,\beta)\varepsilon(\alpha',\beta)$
  and $\varepsilon(\alpha,\beta+\beta')=\varepsilon(\alpha,\beta)\varepsilon(\alpha,\beta')$;
\item $\varepsilon(0,\beta)=\varepsilon(\alpha,0)=1$ and $\varepsilon(\alpha,\beta)^2=1$;
\item the 2-cocycle law \eqref{ch03:eq:coclaw};
\item $\chthreecoc(\alpha,\beta)+\chthreecoc(\beta,\alpha)=\langle\alpha,\beta\rangle
  -\sum_i\langle B_i,B_i\rangle\alpha_i\beta_i$;
\item if $L$ is even, the commutator property \eqref{ch03:eq:commprop}.
\end{enumerate}
\end{theorem}
\begin{leancode}
theorem cocycle_add_left (α α' β : L.carrier) :
    cocycle L B (α + α') β = cocycle L B α β * cocycle L B α' β
theorem cocycle_add_right (α β β' : L.carrier) :
    cocycle L B α (β + β') = cocycle L B α β * cocycle L B α β'
theorem cocycle_zero_left (β : L.carrier) : cocycle L B 0 β = 1
theorem cocycle_zero_right (α : L.carrier) : cocycle L B α 0 = 1
theorem cocycle_sq (α β : L.carrier) : cocycle L B α β * cocycle L B α β = 1
theorem cocycle_cocycleLaw (α β γ : L.carrier) :
    cocycle L B α β * cocycle L B (α + β) γ = cocycle L B β γ * cocycle L B α (β + γ)
theorem cocExp_add_swap (α β : L.carrier) :
    cocExp L B α β + cocExp L B β α
      = L.form α β - ∑ i : Fin n, L.form (B i) (B i) * B.repr α i * B.repr β i
theorem cocycle_comm (hEven : L.Even) (α β : L.carrier) :
    cocycle L B α β * cocycle L B β α = (-1 : ℂ) ^ (L.form α β)
\end{leancode}

\begin{dictionary}
Every statement is universally quantified over an arbitrary basis \lean{B}; the only hypothesis
anywhere is \lean{hEven} in the last one. Since $\varepsilon=\pm1$, the product form of
\eqref{ch03:eq:commprop} is equivalent to the quotient form $\varepsilon(\alpha,\beta)/\varepsilon
(\beta,\alpha)=(-1)^{\langle\alpha,\beta\rangle}$. These are exactly the textbook properties.
\end{dictionary}

\begin{physics}
These identities say that the signs $\varepsilon$ define an associative twisted product on the momentum lattice (the cocycle law) and that $\varepsilon(\alpha,\beta)\varepsilon(\beta,\alpha)=(-1)^{\langle\alpha,\beta\rangle}$: reordering two momentum shifts costs exactly the sign by which the naive vertex operators fail to commute. In CFT terms this is the condition for the Klein-factored vertex operators $:\!e^{i\alpha\cdot X}\!:$ to be mutually local with the correct (bosonic) statistics. The commutator property needs evenness, because otherwise the diagonal term $\sum_i\langle B_i,B_i\rangle\alpha_i\beta_i$ of (iv) does not vanish mod $2$.
\end{physics}

\begin{proof}[Proof on paper]
$\chthreecoc$ is bilinear because coordinates are linear; (i)--(iii) follow by exponentiating.
For (iv), write $\langle\alpha,\beta\rangle=\sum_{i,j}\langle B_i,B_j\rangle\alpha_i\beta_j$ and
split the sum into $j<i$, $j>i$ and $j=i$; by symmetry of the form the $j>i$ part equals
$\chthreecoc(\beta,\alpha)$. For (v), each $\langle B_i,B_i\rangle$ is even, so the diagonal
correction in (iv) is even and disappears from $(-1)^{(\cdot)}$.
\end{proof}

\begin{example}[$A_1$ and $A_2$]\label{ch03:ex:A1A2}
For $A_1=\Z\alpha$ with $\langle\alpha,\alpha\rangle=2$, there is only one basis index, the sum is
empty and $\varepsilon\equiv1$; this is consistent with \eqref{ch03:eq:commprop} because all
pairings $2mn$ are even. For $A_2$ with basis $B_0,B_1$, $\langle B_i,B_i\rangle=2$ and
$\langle B_0,B_1\rangle=-1$, the only pair with $j<i$ is $(i,j)=(1,0)$, so
$\chthreecoc(\alpha,\beta)=-\alpha_1\beta_0$. Thus $\varepsilon(B_1,B_0)=-1$,
$\varepsilon(B_0,B_1)=1$, and the product is $-1=(-1)^{\langle B_0,B_1\rangle}$. Exchanging the
roles of $B_0$ and $B_1$ (another basis) gives a different but cohomologous cocycle; see
Section~\ref{ch03:sec:functoriality}.
\end{example}

\subsection{The canonical cocycle of a lattice}

Since an \lean{IntegralLattice} carries no preferred basis, the library \emph{chooses} one.
\noindent File \leanpath{VOA/Foundations/LatticeCocycle.lean}.
\begin{leancode}
noncomputable def latticeBasis :
    Basis (Fin (Fintype.card (Module.Free.ChooseBasisIndex ℤ L.carrier))) ℤ L.carrier :=
  (Module.Free.chooseBasis ℤ L.carrier).reindex (Fintype.equivFin _)

noncomputable def latticeCocycle (α β : L.carrier) : ℂ := cocycle L (latticeBasis L) α β

theorem latticeCocycle_comm (hEven : L.Even) (α β : L.carrier) :
    latticeCocycle L α β * latticeCocycle L β α = (-1 : ℂ) ^ (L.form α β)
\end{leancode}
together with \lean{latticeCocycle_cocycleLaw}, \lean{latticeCocycle_zero_left},
\lean{latticeCocycle_zero_right}, \lean{latticeCocycle_add_left},
\lean{latticeCocycle_add_right} and \lean{latticeCocycle_sq}, which specialise
Theorem~\ref{ch03:thm:cocycle}.

\begin{dictionary}
\lean{Module.Free.chooseBasis ℤ L.carrier} is \emph{some} basis, extracted from the existence
statement \lean{Module.Free} by the axiom of choice. Lean's \lean{Classical.choice} gives no way to
compute which basis it is, and \lean{reindex (Fintype.equivFin _)} renumbers it by
\lean{Fin n} through another non-constructive bijection. Consequently the value of
\lean{latticeCocycle L α β} for specific $\alpha,\beta$ is a well-defined complex number $\pm1$, but
it \emph{cannot be computed}: no Lean proof can decide by evaluation whether
$\varepsilon(\alpha,\beta)=1$. Every theorem must therefore use only properties valid for
\emph{every} basis. This is a design decision with consequences discussed in
Section~\ref{ch03:sec:design}. It is not a fidelity defect: FLM's construction also depends on a
choice (of central extension), and all choices give isomorphic vertex operator algebras.
\end{dictionary}

\begin{remark}[Further consequences of bimultiplicativity]\label{ch03:rem:negneg}
Because $\varepsilon$ is bimultiplicative and $\pm1$-valued, $\varepsilon(-\alpha,-\beta)=
\varepsilon(\alpha,\beta)^{(-1)(-1)}=\varepsilon(\alpha,\beta)$
(\lean{latticeCocycle_neg_neg}, \leanpath{VOA/Moonshine/LatticeInvolution.lean}), and
$\varepsilon(\alpha,\beta)=(-1)^{\langle\alpha,\beta\rangle}\varepsilon(\beta,\alpha)$
(\lean{latticeCocycle_swap_sign}, \leanpath{VOA/Examples/LatticeGenLocality.lean}):
\begin{leancode}
theorem latticeCocycle_neg_neg (α β : L.carrier) :
    latticeCocycle L (-α) (-β) = latticeCocycle L α β
theorem latticeCocycle_swap_sign (hEven : L.Even) (α β : L.carrier) :
    latticeCocycle L α β = (-1 : ℂ) ^ (L.form α β) * latticeCocycle L β α
\end{leancode}
The first identity is why the lift $\theta$ of $-1\in O(L)$ (Chapter~\ref{ch:04}) needs \emph{no}
sign correction in this library: $\theta(p\otimes e^\alpha)=\sigma(p)\otimes e^{-\alpha}$ exactly.
For a general (non-bimultiplicative) cocycle, FLM must use $\theta e^\alpha=c(\alpha)e^{-\alpha}$
with a correcting function $c$.
\end{remark}

The library also builds the twisted group algebra itself, with product
\lean{twMul} on \lean{L.carrier →₀ ℂ} satisfying \lean{twMul_e}:
$e^\alpha\star e^\beta=\varepsilon(\alpha,\beta)e^{\alpha+\beta}$ and \lean{twMul_assoc}
(\leanpath{VOA/Foundations/LatticeAlgebraTwisted.lean}). For the vertex algebra this product is
never used as such: the cocycle is inserted directly into the vertex operators
(Definition~\ref{ch03:def:Yealpha}), and the state space uses the \emph{vector space}
$\C[L]=$\lean{L.carrier →₀ ℂ}.

\begin{status}
All of Theorem~\ref{ch03:thm:cocycle} and its \lean{latticeCocycle} specialisations: proved,
unconditional except for \lean{hEven} in the commutator property;
\texttt{\#print axioms latticeCocycle\_comm} $=$ \texttt{[propext, Classical.choice, Quot.sound]}
(\lean{Classical.choice} enters, among other places, through \lean{chooseBasis}).
\end{status}

\section{Lattice charges as Heisenberg oscillators: the orthonormal frame}\label{ch03:sec:frame}

The lattice vertex algebra contains the Heisenberg algebra of $\mathfrak h=L\otimes_{\Z}\C$ with the
\emph{lattice} form: $[\alpha(m),\beta(n)]=\langle\alpha,\beta\rangle m\delta_{m+n,0}$. The
library's Fock space $\mathcal F_r$ has \emph{orthonormal} oscillators $h_a(n)$
(Theorem~\ref{ch03:thm:oscN_comm} with $k=1$). To realise $\alpha(n)$ one needs real coordinates
$\alpha^{\perp}\in\chthreeR^r$ with $\sum_a\alpha^\perp_a\beta^\perp_a=\langle\alpha,\beta\rangle$, and then
$\alpha(n):=\sum_a\alpha^\perp_ah_a(n)$. The library obtains such coordinates from the symmetric
positive square root of the Gram matrix.

\begin{definition}[Frame and orthonormal coordinates]\label{ch03:def:frame}
Let $G=(\langle B_i,B_j\rangle)_{ij}$ be the Gram matrix of the chosen basis
\lean{latticeBasis L}, viewed as a real matrix, and $C=\sqrt G$ its positive semidefinite square
root. For $\alpha\in L$ with coordinate vector $\underline\alpha\in\Z^n$, put
$\alpha^\perp=C\underline\alpha\in\chthreeR^n$.
\end{definition}
\begin{leancode}
noncomputable def latRank : ℕ := Fintype.card (Module.Free.ChooseBasisIndex ℤ L.carrier)
noncomputable def gram : Matrix (Fin (latRank L)) (Fin (latRank L)) ℤ :=
  Matrix.of fun i j => L.form (latticeBasis L i) (latticeBasis L j)
noncomputable def coord (α : L.carrier) : Fin (latRank L) → ℤ := fun i => (latticeBasis L).repr α i
noncomputable def gramR : Matrix (Fin (latRank L)) (Fin (latRank L)) ℝ :=
  (gram L).map (Int.cast)
noncomputable def frame : Matrix (Fin (latRank L)) (Fin (latRank L)) ℝ := CFC.sqrt (gramR L)
noncomputable def coordR (α : L.carrier) : Fin (latRank L) → ℝ := fun i => ((coord L α i : ℤ) : ℝ)
noncomputable def orthCoord (α : L.carrier) : Fin (latRank L) → ℝ := frame L *ᵥ coordR L α

theorem orthCoord_inner (hPos : L.PosDef) (α β : L.carrier) :
    ∑ a, orthCoord L α a * orthCoord L β a = ((L.form α β : ℤ) : ℝ)
\end{leancode}
(\leanpath{VOA/Foundations/LatticeFinite.lean}, \leanpath{VOA/Foundations/LatticeOrthoFrame.lean}.)

\begin{dictionary}
\lean{CFC.sqrt} is Mathlib's continuous-functional-calculus square root; on a positive
semidefinite real symmetric matrix it is the unique positive semidefinite square root.
\lean{*ᵥ} is matrix--vector multiplication. \lean{orthCoord_inner} is the identity
$\alpha^\perp\cdot\beta^\perp=\underline\alpha^TC^TC\underline\beta=\underline\alpha^TG\underline\beta
=\langle\alpha,\beta\rangle$, valid because $C$ is symmetric and $C^2=G$; it needs \lean{hPos}
only to know $G$ is positive semidefinite so that the square root behaves. The integers
\lean{latRank L} and \lean{rank L} are equal (\lean{rank_eq_latRank}), and
\lean{rankEquiv : Fin (rank L) ≃ Fin (latRank L)} transports indices between the Fock colours and
the frame coordinates.
\end{dictionary}

\begin{physics}
The frame $\alpha\mapsto\alpha^\perp$ realises the abstract lattice as a lattice of momentum vectors in the Euclidean space $\mathbb R^r$ with orthonormal coordinates: the $a$-th component $\alpha^\perp_a$ is the momentum carried by the $a$-th boson, and $\langle\alpha,\beta\rangle=\alpha^\perp\cdot\beta^\perp$. Choosing the frame is choosing Cartesian coordinates on the target torus, which is how one passes from $r$ independent bosons to the vertex operator $:\!e^{i\alpha^\perp\cdot X}\!:$. It is pure bookkeeping (a change of basis), with no new physics, but all later formulas are written in these coordinates.
\end{physics}

\begin{definition}[Charge oscillators]\label{ch03:def:chargeOsc}
$\alpha(m):=\sum_{a}\alpha^\perp_a\,h_a(m)$ on $\mathcal F_{\operatorname{rank}L}$ (level $k=1$).
\end{definition}
\begin{leancode}
noncomputable def chargeOsc (α : L.carrier) (m : ℤ) :
    Module.End ℂ (HeisenbergN.FockN (rank L)) :=
  ∑ a : Fin (rank L), (orthCoord L α (rankEquiv L a) : ℂ) • HeisenbergN.oscN (rank L) 1 a m

noncomputable def chargeCoord (α : L.carrier) (b : Fin (rank L)) : ℂ :=
  (orthCoord L α (rankEquiv L b) : ℂ)

theorem chargeOsc_comm (hPos : L.PosDef) (α β : L.carrier) (m n : ℤ) :
    chargeOsc L α m * chargeOsc L β n - chargeOsc L β n * chargeOsc L α m
      = (if m + n = 0 then ((L.form α β : ℤ) : ℂ) * (m : ℂ) else 0) • LinearMap.id
\end{leancode}
(\leanpath{VOA/Examples/LatticeGenCharge.lean}, \leanpath{LatticeGenChargeFields.lean}.)

\begin{physics}
The charge oscillator $\alpha(m)=\alpha^\perp\cdot h(m)$ is the mode of the current component $\alpha\cdot i\partial X(z)$ along the direction $\alpha$ in target space, and its zero mode $\alpha(0)$ measures the momentum of a state projected onto $\alpha$. Its commutator $[\alpha(m),\beta(n)]=m\langle\alpha,\beta\rangle\delta_{m+n,0}$ is the OPE of two such projected currents with the lattice inner product as the two-point coefficient. It is the building block of the vertex operator $:\!e^{i\alpha\cdot X}\!:$, whose exponent is the integral of this current.
\end{physics}

\begin{proof}[Proof on paper of \lean{chargeOsc_comm}]
Expand both oscillators in the $h_a$; by Theorem~\ref{ch03:thm:oscN_comm} only the diagonal
$a=b$ terms survive and give $m\delta_{m+n,0}\sum_a\alpha_a^\perp\beta_a^\perp=m\delta_{m+n,0}
\langle\alpha,\beta\rangle$.
\end{proof}

\begin{remark}[Why $\sqrt G$, and what it costs]\label{ch03:rem:sqrtG}
Any real matrix $C$ with $C^TC=G$ would do; $\sqrt G$ is chosen because it exists canonically in
Mathlib. The price is that the oscillator directions $h_a$ of the library are \emph{not} adapted
to any interesting coordinate system. For the Leech lattice, for instance, the $24$ Fock colours
are not the coordinate directions of the Golay-code construction. A state such as
$h_a(-1)\mathbb 1$ (\lean{heisSt}) is an orthonormal-frame state with no lattice meaning; the
lattice-meaningful states are $\alpha(-1)\mathbb 1$ (\lean{hSt}). This is the second half of the
design question of Section~\ref{ch03:sec:design}.
\end{remark}

\section{The state space of \texorpdfstring{$V_L$}{V\_L} and its grading}\label{ch03:sec:latticespace}

\subsection{The underlying vector space}

\begin{definition}\label{ch03:def:LatModuleGen}
$V_L:=\mathcal F_{\operatorname{rank}L}\otimes_{\C}\C[L]$, where $\C[L]$ is the vector space with
basis $\{e^\alpha\}_{\alpha\in L}$. Write $p\otimes e^\alpha$ for basis-like states, $\mathbb 1=
1\otimes e^0$ for the vacuum and $e^\alpha=1\otimes e^\alpha$ for the charged vacua.
\end{definition}
\begin{leancode}
noncomputable abbrev LatModuleGen : Type :=
  HeisenbergN.FockN (rank L) ⊗[ℂ] (L.carrier →₀ ℂ)

noncomputable def eGen (α : L.carrier) : L.carrier →₀ ℂ := Finsupp.single α 1

noncomputable def latStateGen (p : HeisenbergN.FockN (rank L)) (α : L.carrier) : LatModuleGen L :=
  p ⊗ₜ[ℂ] eGen L α

noncomputable def latVacuumGen : LatModuleGen L :=
  latStateGen L (HeisenbergN.vacuumN (rank L)) 0

noncomputable def latChargeGen (α : L.carrier) : LatModuleGen L :=
  latStateGen L (HeisenbergN.vacuumN (rank L)) α
\end{leancode}
(\leanpath{VOA/Examples/LatticeGenModule.lean}.)

\begin{dictionary}
\lean{L.carrier →₀ ℂ} is the space of finitely supported functions $L\to\C$, i.e.\ the vector
space with basis $e^\alpha=$ \lean{Finsupp.single α 1}. \lean{⊗[ℂ]} is the algebraic tensor
product over $\C$ and \lean{p ⊗ₜ[ℂ] x} a pure tensor. The cocycle is \emph{not} part of the
state space; it appears only in the vertex operators. The states $p\otimes e^\alpha$ span $V_L$
(\lean{latStateGen_span}).
\end{dictionary}

\begin{physics}
The space $V_L=\mathcal F_r\otimes\C[L]$ is the full space of states of $r$ compactified chiral bosons: the Fock factor holds the oscillator excitations and the group algebra $\C[L]$ holds the discrete momentum $\alpha\in L$, with $e^\alpha$ the ground state of momentum $\alpha$. The charged vacuum $e^\alpha$ is the state created from the vacuum by the vertex operator $:\!e^{i\alpha\cdot X}\!:$. Decomposing $V_L=\bigoplus_\gamma\mathcal F\otimes e^\gamma$ is the decomposition by $U(1)^r$ charge (the momentum sectors), and the chosen $\gamma$ sector is the charge-$\gamma$ module of the Heisenberg algebra.
\end{physics}

We also need the inclusion of the charge-$\gamma$ sector and the extraction of the
charge-$\gamma$ component:
\noindent File \leanpath{VOA/Examples/LatticeCharge.lean}.
\begin{leancode}
noncomputable def mkEα (α : L.carrier) : HeisenbergN.FockN (rank L) →ₗ[ℂ] LatModuleGen L :=
  (TensorProduct.mk ℂ (HeisenbergN.FockN (rank L)) (L.carrier →₀ ℂ)).flip (eGen L α)

noncomputable def coordEα (α : L.carrier) : LatModuleGen L →ₗ[ℂ] HeisenbergN.FockN (rank L) :=
  (TensorProduct.rid ℂ (HeisenbergN.FockN (rank L))).toLinearMap.comp
    (LinearMap.lTensor (HeisenbergN.FockN (rank L)) (Finsupp.lapply α))

noncomputable def chargeSec (γ : L.carrier) : Submodule ℂ (LatModuleGen L) :=
  LinearMap.range (mkEα L γ)
\end{leancode}
\lean{mkEα L α p} $=p\otimes e^\alpha$; \lean{coordEα L γ} sends $p\otimes e^{\alpha}$ to $p$ if
$\alpha=\gamma$ and to $0$ otherwise (\lean{coordEα_apply}), so that every $x\in V_L$ is
$\sum_\gamma\lean{coordEα}_\gamma(x)\otimes e^\gamma$, a finite sum; and \lean{chargeSec L γ}
$=\mathcal F\otimes e^\gamma$ is the \emph{momentum-$\gamma$ sector}.

\subsection{The conformal weight grading}

The conformal weight of $p\otimes e^\alpha$, with $p$ of Fock weight $w$, is
$w+\tfrac12\langle\alpha,\alpha\rangle$.

\begin{definition}[Grading of $V_L$]\label{ch03:def:gradingLatGen}
For $\alpha\in L$ let $|\alpha|^2/2:=\langle\alpha,\alpha\rangle/2$ (an integer when $L$ is even).
For $(w,\alpha)\in\Z\times L$ let $V_{(w,\alpha)}=(\mathcal F)_w\otimes e^\alpha$ (zero if $w<0$),
and for $n\in\Z$ let $(V_L)_n=\sum_\alpha V_{(n-|\alpha|^2/2,\alpha)}$.
\end{definition}
\begin{leancode}
noncomputable def halfNorm (α : L.carrier) : ℤ := L.form α α / 2

noncomputable def bigradα (wα : ℤ × L.carrier) : Submodule ℂ (LatModuleGen L) :=
  (HeisenbergN.gradingZN (rank L) wα.1).map (mkEα L wα.2)

noncomputable def gradingLatGen (k : ℤ) : Submodule ℂ (LatModuleGen L) :=
  ⨆ α : L.carrier, bigradα L (k - halfNorm L α, α)
\end{leancode}
(\leanpath{VOA/Examples/LatticeGenGrading.lean}.)

\begin{dictionary}
\lean{L.form α α / 2} is integer division; it equals $\langle\alpha,\alpha\rangle/2$ exactly when
$\langle\alpha,\alpha\rangle$ is even, which is the content of
\begin{leancode}
theorem two_halfNorm (hEven : L.Even) (α : L.carrier) : 2 * halfNorm L α = L.form α α
\end{leancode}
i.e.\ $2\cdot$\lean{halfNorm L α} $=\langle\alpha,\alpha\rangle$. Without evenness \lean{halfNorm} would
silently round, but every theorem using it assumes \lean{hEven}. \lean{Submodule.map} is the
image of a submodule under a linear map, and \lean{⨆ α} is the supremum (sum) of a family of
submodules. So \lean{gradingLatGen L k} is literally $\sum_\alpha(\mathcal F)_{k-|\alpha|^2/2}
\otimes e^\alpha$.
\end{dictionary}

\begin{physics}
This is the conformal-weight formula for a state of momentum $\alpha$ and oscillator level $w$: $h=\frac12\langle\alpha,\alpha\rangle+w$. The first term is the kinetic energy $p^2/2$ of the zero mode, the second the energy of the excited oscillators. The grading is the eigenvalue of $L_0$ on the cylinder (up to the $-c/24$ shift), and evenness makes it integral. The weight-one subspace $(V_L)_1$ is then the space of conserved currents: oscillators $h_a(-1)\mathbb1$ plus the momentum states of norm $2$ (the roots).
\end{physics}

\begin{theorem}[The grading is a graded decomposition of CFT type]\label{ch03:thm:gradingLat}
Let $L$ be even and positive-definite. Then $V_L=\bigoplus_n(V_L)_n$, each $(V_L)_n$ is
finite-dimensional, $(V_L)_n=0$ for $n<0$ and $(V_L)_0=\C\mathbb 1$.
\end{theorem}
\begin{leancode}
theorem gradingLatGen_isInternal_of_posDef (hEven : L.Even) (hPos : L.PosDef) :
    DirectSum.IsInternal (gradingLatGen L)
theorem gradingLatGen_finrank_of_posDef (hEven : L.Even) (hPos : L.PosDef) (k : ℤ) :
    FiniteDimensional ℂ (gradingLatGen L k)
theorem gradingLatGen_positivity (hEven : L.Even) (hPos : L.PosDef) (k : ℤ) (hk : k < 0) :
    gradingLatGen L k = ⊥
theorem gradingLatGen_zero_eq_span_vacuum (hEven : L.Even) (hPos : L.PosDef) :
    gradingLatGen L 0 = ℂ ∙ latVacuumGen L
\end{leancode}

\begin{physics}
This is the statement that the lattice theory has a well-defined, discrete energy spectrum with a unique vacuum: each energy level has finitely many states (because only finitely many lattice vectors have bounded norm, by positive-definiteness, and each Fock level is finite), no negative energies occur (unitarity-type positivity), and only the vacuum has $h=0$ (no zero-weight vertex operator other than the identity). These are the conditions for $V_L$ to be of CFT type, the algebraic version of a unitary chiral CFT with a normalisable vacuum.
\end{physics}
\begin{proof}[Proof on paper]
The bigrading $V_L=\bigoplus_{(w,\alpha)}(\mathcal F)_w\otimes e^\alpha$ is a direct sum because
$\mathcal F=\bigoplus_w\mathcal F_w$ and $\C[L]=\bigoplus_\alpha\C e^\alpha$; collapsing along the
fibres of $(w,\alpha)\mapsto w+|\alpha|^2/2$ preserves direct sums. If $n<0$ then
$n-|\alpha|^2/2<0$ for every $\alpha$ (positivity of the norm), so $(V_L)_n=0$. The weight-$0$
space needs $w=0$ and $|\alpha|^2=0$, hence $\alpha=0$ and the state is a multiple of the vacuum.
Finite-dimensionality: only the finitely many $\alpha$ with $|\alpha|^2\le2n$ contribute
(\lean{boundedFinite}), and each contributes the finite-dimensional $\mathcal F_{n-|\alpha|^2/2}$.
\end{proof}
\begin{status}
Proved, conditional only on \lean{hEven} and \lean{hPos}. (The file's module docstring still
calls the finiteness of bounded vectors an ``isolated crux''; that crux has since been
discharged by \lean{boundedFinite}, and the \lean{\_of\_posDef} versions above use it. The
docstring is stale; the theorems are not.)
\end{status}

Consequently the graded dimension is
\begin{equation}\label{ch03:eq:latticechar}
  \sum_{n}\dim(V_L)_n\,q^n=\frac{\theta_L(q)}{\prod_{n\geq1}(1-q^n)^{r}},\qquad
  \theta_L(q)=\sum_{\alpha\in L}q^{\langle\alpha,\alpha\rangle/2},
\end{equation}
which for the Leech lattice is the subject of Chapter~\ref{ch:05}.

\section{The vertex operators \texorpdfstring{$Y(e^\alpha,z)$}{Y(e\^{}alpha,z)}}\label{ch03:sec:vertexops}

\subsection{The exponentials \texorpdfstring{$E^\pm$}{E+-}}

The two exponential factors in \eqref{ch03:eq:VOphysical} are
\begin{align}
  E^-(\alpha,z)&=\exp\Bigl(\sum_{n\ge1}\frac{\alpha(-n)}{n}z^n\Bigr)=\sum_{N\ge0}S_N(\alpha)z^N,
  \label{ch03:eq:Epm}\\
  E^+(\alpha,z)&=\exp\Bigl(-\sum_{n\ge1}\frac{\alpha(n)}{n}z^{-n}\Bigr)=\sum_{N\ge0}E^+_N(\alpha)z^{-N}.
  \notag
\end{align}
\emph{Naming caution}: these are FLM's $E^-(-\alpha,z)$ and $E^+(-\alpha,z)$; the library drops
the minus sign from the name. The coefficients $S_N(\alpha)$ are the elementary Schur polynomials
in the creation operators, and the library computes them, without any exponential, by the
Newton-type recursion obtained by differentiating: if $E(z)=\exp(\sum_{n\ge1} a_nz^n/n)=
\sum_NS_Nz^N$ (the $a_n$ commuting), then $E'(z)=(\sum_{n\ge1}a_nz^{n-1})E(z)$; comparing the
coefficients of $z^N$ on both sides gives $(N+1)S_{N+1}=\sum_{m=0}^{N}a_{m+1}S_{N-m}$, which
determines all $S_N$ from $S_0=1$. Conversely, the series defined by the recursion satisfies the
same first-order differential equation with the same constant term, hence equals $E(z)$.

\begin{definition}[Schur coefficients]\label{ch03:def:vSchur}
$v_m(\alpha):=\sum_a\alpha^\perp_aX_{(m,a)}$ is the state $\alpha(-m-1)\mathbb 1$. Define
$S_0(\alpha)=1$, $S_{N+1}(\alpha)=\frac1{N+1}\sum_{m=0}^{N}v_m(\alpha)S_{N-m}(\alpha)$, and the
annihilation coefficients $E^+_0=\mathrm{id}$,
$E^+_{N+1}=\frac{1}{N+1}\sum_{m=0}^N(-\alpha(m+1))\circ E^+_{N-m}$.
\end{definition}
\begin{leancode}
noncomputable def vVar (α : L.carrier) (m : ℕ) : HeisenbergN.FockN (rank L) :=
  ∑ a : Fin (rank L), chargeCoord L α a • (MvPolynomial.X ((m, a)) : HeisenbergN.FockN (rank L))

noncomputable def vSchur (α : L.carrier) : ℕ → HeisenbergN.FockN (rank L)
  | 0 => 1
  | (N+1) => ((N+1 : ℂ))⁻¹ • ∑ m ∈ Finset.range (N+1), vVar L α m * vSchur α (N - m)

noncomputable def EminusVec (α : L.carrier) (N : ℕ) : End ℂ (HeisenbergN.FockN (rank L)) :=
  LinearMap.mulLeft ℂ (vSchur L α N)

noncomputable def aExpVec (α : L.carrier) : ℕ → End ℂ (HeisenbergN.FockN (rank L))
  | 0 => LinearMap.id
  | (N+1) => ((N+1 : ℂ))⁻¹ • ∑ m ∈ Finset.range (N+1),
      (-(chargeOsc L α (m+1 : ℤ))).comp (aExpVec α (N - m))
\end{leancode}
(\leanpath{VOA/Examples/LatticeGenChargeFields.lean}.) Since creation operators are
multiplications, $E^-$ acts by multiplication by the polynomial $S_N(\alpha)$
(\lean{LinearMap.mulLeft}). As formal distributions (indexed by exponents of $z$):
\begin{leancode}
noncomputable def EminusVecField (α : L.carrier) : FormalDist (End ℂ (FockN (rank L))) :=
  fun t => if 0 ≤ t then EminusVec L α t.toNat else 0

noncomputable def EplusVecField (α : L.carrier) : FormalDist (End ℂ (FockN (rank L))) :=
  fun t => if t ≤ 0 then aExpVec L α (-t).toNat else 0

noncomputable def EEVec (α : L.carrier) : FormalDist (End ℂ (FockN (rank L))) :=
  nprod (EminusVecField_isField L α) (EplusVecField_isField L α)
\end{leancode}
(\leanpath{VOA/Examples/LatticeGenVertex.lean}.)

\begin{dictionary}
\lean{EminusVecField} has nonzero coefficients only at $z^t$, $t\ge0$, and there it is
$S_t(\alpha)$; \lean{EplusVecField} only at $z^t$, $t\le0$, and there it is $E^+_{-t}(\alpha)$.
\lean{EEVec} is the normal-ordered product ${:}E^-(\alpha,z)E^+(\alpha,z){:}$; since all
nonzero coefficients of $E^-$ sit at non-negative powers, the normal ordering simply puts $E^-$ on
the left, so \lean{EEVec} $=E^-(\alpha,z)E^+(\alpha,z)$, the boson dressing. It acts on the Fock
factor only.
\end{dictionary}

\begin{physics}
The coefficients $S_N(\alpha)$ are the Schur polynomials in the creation modes and arise from expanding the exponential of the negative part of the current, $\exp\bigl(\sum_{m>0}\frac{\alpha(-m)}{m}z^m\bigr)=\sum_NS_N(\alpha)z^N$. They are the oscillator content of the vertex operator $:\!e^{i\alpha\cdot X}\!:$ in its factorised form (the creation part $E^-$ and the annihilation part $E^+$), obtained by the standard normal-ordering of the exponential of a free field. This is technical infrastructure for writing $Y(e^\alpha,z)$ explicitly.
\end{physics}

\subsection{The charged vertex operator}

Before the definition, one convention must be fixed: in $Y(e^\alpha,z)=E^-E^+e_\alpha
z^{\alpha(0)}$ the operator $z^{\alpha(0)}$ acts \emph{first}, on the incoming momentum $\beta$,
giving $z^{\langle\alpha,\beta\rangle}$, and then $e_\alpha$ shifts $\beta$ to $\alpha+\beta$ with
the sign $\varepsilon(\alpha,\beta)$. This is FLM's convention (8.4.1).

\begin{definition}[$Y(e^\alpha,z)$]\label{ch03:def:Yealpha}
For $\alpha\in L$, define $Y(e^\alpha,z)$ on $p\otimes e^\beta$ by
\begin{equation}\label{ch03:eq:Yealpha}
  Y(e^\alpha,z)(p\otimes e^\beta)=\varepsilon(\alpha,\beta)\,z^{\langle\alpha,\beta\rangle}\,
  \bigl(E^-(\alpha,z)E^+(\alpha,z)p\bigr)\otimes e^{\alpha+\beta}.
\end{equation}
Equivalently, its $z^t$-coefficient sends $p\otimes e^\beta$ to
$\varepsilon(\alpha,\beta)\,[E^-E^+]_{t-\langle\alpha,\beta\rangle}\,p\otimes e^{\alpha+\beta}$.
\end{definition}
\begin{leancode}
noncomputable def YeαvecFam (α : L.carrier) (t : ℤ) (β : L.carrier) :
    ℂ →ₗ[ℂ] (FockN (rank L) →ₗ[ℂ] LatModuleGen L) :=
  LinearMap.toSpanSingleton ℂ (FockN (rank L) →ₗ[ℂ] LatModuleGen L)
    (latticeCocycle L α β •
      (((TensorProduct.mk ℂ (FockN (rank L)) (L.carrier →₀ ℂ)).flip (eGen L (α + β))).comp
        (EEVec L α (t - L.form α β))))

noncomputable def Yeαvec (α : L.carrier) (t : ℤ) : LatModuleGen L →ₗ[ℂ] LatModuleGen L :=
  TensorProduct.lift (Finsupp.lsum ℂ (YeαvecFam L α t)).flip

theorem Yeαvec_tmul (α β : L.carrier) (t : ℤ) (p : FockN (rank L)) :
    Yeαvec L α t (latStateGen L p β)
      = latticeCocycle L α β • latStateGen L (EEVec L α (t - L.form α β) p) (α + β)
\end{leancode}

\begin{dictionary}
The first two declarations are the bureaucracy of defining a linear map on a tensor product with
a finitely supported factor; the third, \lean{Yeαvec_tmul}, is the defining formula and is exactly
the coefficient form of \eqref{ch03:eq:Yealpha}: the factor $z^{\alpha(0)}$ acts on the
$\beta$-sector as $z^{\langle\alpha,\beta\rangle}$, which shifts the exponent \lean{t} of the
dressing by $\langle\alpha,\beta\rangle$; $e_\alpha$ shifts the momentum to $\alpha+\beta$ and
contributes $\varepsilon(\alpha,\beta)$. Note the order convention $e_\alpha z^{\alpha(0)}$ with
$z^{\alpha(0)}$ evaluated on the \emph{incoming} momentum $\beta$; this is FLM's convention
(8.4.1). With it, $Y(e^\alpha,z)\mathbb 1=E^-(\alpha,z)e^\alpha$, so $(e^\alpha)_{(-1)}\mathbb 1=
e^\alpha$ (\lean{Yeαvec_mode_neg_one_vacuum}), and
\begin{equation*}
  (e^\alpha)_{(-\langle\alpha,\beta\rangle-1)}e^\beta=\varepsilon(\alpha,\beta)e^{\alpha+\beta}
\end{equation*}
(\lean{smul_latChargeGen}, used in Section~\ref{ch03:sec:rigidity}).
\end{dictionary}

\begin{physics}
This is the vertex operator $:\!e^{i\alpha\cdot X(z)}\!:=E^-(\alpha,z)E^+(\alpha,z)\,e_\alpha\,z^{\alpha(0)}$ of the compactified boson: it shifts the momentum by $\alpha$, the factor $z^{\alpha(0)}$ produces the OPE singularity $z^{\langle\alpha,\beta\rangle}$ with a state of momentum $\beta$, the exponentials $E^\mp$ are its oscillator dressing, and the sign $\varepsilon(\alpha,\beta)$ is the Klein factor. Its conformal weight is $\langle\alpha,\alpha\rangle/2$, and acting on the vacuum it creates the charged ground state $e^\alpha$. This is the central formula of Frenkel--Kac and Segal (the construction of simply-laced affine algebras from lattices) and of FLM Chapter 8.
\end{physics}

The boson currents on $V_L$ are the \emph{lifted} currents (the Fock currents acting on the first
tensor factor only, $h_a(n)\otimes\mathrm{id}$) plus the momentum zero mode:
$h_a(z)=\sum_nh_a(n)z^{-n-1}$ with $h_a(0)(p\otimes e^\beta)=\beta^\perp_a\,p\otimes e^\beta$.
\begin{leancode}
noncomputable def zGradeOpGen (a : Fin (rank L)) : LatModuleGen L →ₗ[ℂ] LatModuleGen L :=
  LinearMap.lTensor (FockN (rank L)) (zGradeLGen L a)

theorem zGradeOpGen_latStateGen (a : Fin (rank L)) (p : FockN (rank L)) (β : L.carrier) :
    zGradeOpGen L a (latStateGen L p β)
      = (chargeCoord L β a) • latStateGen L p β

noncomputable def alphaLatGen (a : Fin (rank L)) : FormalDist (End ℂ (LatModuleGen L)) :=
  fun j => LinearMap.rTensor (L.carrier →₀ ℂ) (currentN (rank L) 1 a j)
    + (if j = -1 then zGradeOpGen L a else 0)

noncomputable def genFieldLatGen (a : Fin (rank L)) (i : ℕ) : FormalDist (End ℂ (LatModuleGen L)) :=
  (i.factorial : ℂ)⁻¹ • (VOA.deriv^[i] (alphaLatGen L a))
\end{leancode}
(\leanpath{VOA/Examples/LatticeGenBosonFields.lean}.) The coefficient of $z^{-1}$ is the mode
$h_a(0)$, which is why the zero-mode operator is added at \lean{j = -1}. \lean{rTensor f} is
$f\otimes\mathrm{id}$, \lean{lTensor g} is $\mathrm{id}\otimes g$, and \lean{VOA.deriv^[i]} is the
$i$-fold formal derivative.

\section{Locality of the lattice vertex operators}\label{ch03:sec:locality}

\begin{theorem}[Charge-pair locality]\label{ch03:thm:Yelocal}
Let $L$ be even and positive-definite. For all $\alpha,\beta\in L$, $Y(e^\alpha,z)$ and
$Y(e^\beta,z)$ are mutually local.
\end{theorem}
\begin{leancode}
theorem Yeαvec_local (hEven : L.Even) (hPos : L.PosDef) (α β : L.carrier) :
    Local (Yeαvec L α) (Yeαvec L β)
\end{leancode}
(\leanpath{VOA/Examples/LatticeGenOppLocal.lean}.)

\begin{dictionary}
\lean{Local} is Definition-level locality of Section~\ref{ch03:sec:conventions}: some power
$(z_1-z_2)^N$ annihilates the commutator. The hypothesis \lean{hEven} is essential (it gives the
commutator property of $\varepsilon$); \lean{hPos} enters through the frame
(\lean{orthCoord_inner}), which needs a positive semidefinite Gram matrix to take its square
root. The theorem is not vacuous: $Y(e^\alpha,z)$ is a nonzero field for every $\alpha$ (it creates
$e^\alpha$ from the vacuum).
\end{dictionary}

\begin{physics}
Mutual locality of two vertex operators is the statement that their correlators are meromorphic with only poles at coincident points, so that the operators commute (or anticommute for fermions) at separated points: $(z-w)^N[Y(e^\alpha,z),Y(e^\beta,w)]=0$. Here the OPE gives the factor $(z-w)^{\langle\alpha,\beta\rangle}$ times a regular function, which is single-valued because $\langle\alpha,\beta\rangle\in\Z$; the Klein factors $\varepsilon$ cancel the sign $(-1)^{\langle\alpha,\beta\rangle}$ from exchanging the operators so that $e^\alpha$ and $e^\beta$ are genuinely bosonic. This is the key compatibility condition that makes the lattice theory a vertex algebra rather than a mere collection of operators.
\end{physics}

\begin{proof}[Proof on paper]
The one computation behind everything is the commutation of the two exponentials. Since
$[\beta(n),\alpha(-m)]=n\langle\alpha,\beta\rangle\delta_{nm}$ is a scalar multiple of the identity,
the identity $e^Ae^B=e^{[A,B]}e^Be^A$ (valid whenever $[A,B]$ is a scalar) gives
\begin{equation*}
 E^+(\beta,z_2)E^-(\alpha,z_1)=\exp\Bigl(-\sum_{n\ge1}\frac{\langle\alpha,\beta\rangle}{n}
 \Bigl(\frac{z_1}{z_2}\Bigr)^n\Bigr)E^-(\alpha,z_1)E^+(\beta,z_2)
 =\Bigl(1-\frac{z_1}{z_2}\Bigr)^{\langle\alpha,\beta\rangle}E^-(\alpha,z_1)E^+(\beta,z_2),
\end{equation*}
expanded in non-negative powers of $z_1/z_2$. Put $P=\langle\alpha,\beta\rangle$ and act on
$p\otimes e^\gamma$. Using bimultiplicativity of $\varepsilon$,
\begin{align*}
 Y(e^\alpha,z_1)Y(e^\beta,z_2)(p\otimes e^\gamma)&=\varepsilon(\alpha,\beta)\varepsilon(\alpha,\gamma)
 \varepsilon(\beta,\gamma)\,z_1^{\langle\alpha,\gamma\rangle}z_2^{\langle\beta,\gamma\rangle}
 \,(z_1-z_2)^P\;\mathcal N\,p\otimes e^{\alpha+\beta+\gamma},\\
 Y(e^\beta,z_2)Y(e^\alpha,z_1)(p\otimes e^\gamma)&=\varepsilon(\beta,\alpha)\varepsilon(\alpha,\gamma)
 \varepsilon(\beta,\gamma)\,z_1^{\langle\alpha,\gamma\rangle}z_2^{\langle\beta,\gamma\rangle}
 \,(z_2-z_1)^P\;\mathcal N\,p\otimes e^{\alpha+\beta+\gamma},
\end{align*}
where $\mathcal N=E^-(\alpha,z_1)E^-(\beta,z_2)E^+(\alpha,z_1)E^+(\beta,z_2)$ is symmetric under the
exchange, $(z_1-z_2)^P$ is expanded in $|z_1|>|z_2|$ and $(z_2-z_1)^P$ in $|z_2|>|z_1|$. By the
commutator property $\varepsilon(\alpha,\beta)=(-1)^P\varepsilon(\beta,\alpha)$, and
$(z_2-z_1)^P=(-1)^P(z_1-z_2)^P$ formally. Two cases:
\begin{itemize}
\item $P\geq0$: $(z_1-z_2)^P$ is a polynomial, both expansions agree, the two signs cancel and the
  commutator is \emph{exactly zero} ($N=0$).
\item $P<0$: the two expansions of $(z_1-z_2)^P$ differ, but after multiplication by
  $(z_1-z_2)^{-P}$ both become $1$; hence $(z_1-z_2)^{-P}[Y(e^\alpha,z_1),Y(e^\beta,z_2)]=0$.
\end{itemize}
In the second case one also needs that multiplying the formal expansions by
$(z_1-z_2)^{-P}$ interacts correctly with the infinite sums, which is a finite bookkeeping on each
vector because $Y(e^\alpha,z)$ is a field (\lean{Yeαvec_isField}). The two cases are the files
\leanpath{LatticeGenSameSign*.lean} and \leanpath{LatticeGenOpp*.lean}.
\end{proof}

\begin{status}
Proved, conditional on \lean{hEven}, \lean{hPos}; part of the closure of
\lean{latYfullGen_of_locality}, whose axioms are the standard three.
\end{status}

\begin{warning}
Without the cocycle the same computation gives $[Y(e^\alpha,z_1),Y(e^\beta,z_2)]\ne0$ whenever
$P$ is odd and $\geq 0$: the fields \emph{anticommute}. The rank-$1$ lattice $A_1$ was built
earlier in the library with $\varepsilon\equiv1$ (\leanpath{VOA/Examples/LatticeA1*.lean}); this is
correct only because all pairings in $A_1$ are even. The general construction must not be
specialised to $\varepsilon\equiv 1$ for, e.g., $A_2$ or the Leech lattice, where odd pairings
occur.
\end{warning}

The currents are local with each other ($N=2$) and with each $Y(e^\beta,z)$
(\lean{genFieldLatGen_local_Yeαvec}); the latter is the statement
$[h_a(m),Y(e^\beta,z)]=\beta^\perp_az^mY(e^\beta,z)$ in disguise, i.e.\ $(z_1-z_2)$ kills the
commutator.

\section{Reconstruction from locality and the instance \texorpdfstring{$V_L$}{V\_L}}\label{ch03:sec:reconstruction}

\subsection{The strategy}

The textbook definition of $Y(v,z)$ for a general state $v=h_{a_1}(-i_1-1)\cdots
h_{a_s}(-i_s-1)e^\alpha$ is
\begin{equation}\label{ch03:eq:Ygeneral}
  Y(v,z)={:}\partial^{(i_1)}h_{a_1}(z)\cdots\partial^{(i_s)}h_{a_s}(z)\,Y(e^\alpha,z){:},
\end{equation}
after which one must check all vertex algebra axioms, notably the Jacobi identity or locality for
\emph{every} pair of states. The library instead uses the \emph{reconstruction theorem} (Kac's
``existence theorem''): it is enough to have a family of mutually local fields whose creation
states span, together with a translation operator $T$ for which the fields are covariant. Then a
unique vertex map exists on the whole space and satisfies creation, locality and covariance.

\begin{definition}[Translation covariance]\label{ch03:def:DCovariant}
A formal distribution $a$ is $T$-covariant if $[T,a(z)]=\partial_za(z)$. (This is the
infinitesimal form of translation invariance $e^{wT}a(z)e^{-wT}=a(z+w)$; in a vertex operator
algebra $T=L_{-1}$.) In the $z$-exponent
indexing: $[T,a_m]=(m+1)a_{m+1}$ for the $z^m$-coefficient $a_m$.
\end{definition}
\noindent File \leanpath{VOA/Core/ReconstructionExistence.lean}.
\begin{leancode}
def DCovariant (T : End ℂ V) (a : FormalDist (End ℂ V)) : Prop :=
  ∀ m : ℤ, T * a m - a m * T = ((m + 1 : ℤ)) • a (m + 1)
\end{leancode}

\begin{physics}
The operator $T=L_{-1}$ generates translations in the plane coordinate $z$, and $[T,a(z)]=\partial_za(z)$ is the statement that $a(z)$ is a local field, whose position dependence is generated by $T$. In the radial picture $e^{wT}a(z)e^{-wT}=a(z+w)$ is just translation covariance of a quantum field. In a CFT this holds for every field automatically once $T=L_{-1}$ is the stress-tensor mode; here it is a hypothesis of the reconstruction theorem, imposed before a conformal vector is available.
\end{physics}

\begin{theorem}[Reconstruction]\label{ch03:thm:reconstruct}
Let $V$ be a vector space with a vector $\mathbb 1$ and $T\in\End V$ with $T\mathbb1=0$. Let
$(a^i)_{i\in I}$ be fields on $V$ and $(s_i)_{i\in I}$ vectors spanning $V$, such that the $a^i$
are pairwise local, $T$-covariant, and $a^i(z)\mathbb 1=s_i+O(z)$ (i.e.\ the $z^0$-coefficient on
the vacuum is $s_i$ and negative powers vanish). Then there is a linear map
$Y\colon V\to\End(V)[[z^{\pm1}]]$ with $Y(s_i)=a^i$ for all $i$, every $Y(u)$ a field,
$Y(u),Y(v)$ local, $Y(u,z)\mathbb1=u+O(z)$ and $Y(u)$ $T$-covariant, for all $u,v$.
\end{theorem}
\begin{leancode}
theorem reconstruct {ι : Type*} (vac : V) (st : ι → V) (fld : ι → FormalDist (End ℂ V))
    (T : End ℂ V)
    (hspan : Submodule.span ℂ (Set.range st) = ⊤)
    (hfield : ∀ i, IsField (fld i))
    (hlocal : ∀ i j, Local (fld i) (fld j))
    (hcre0 : ∀ i, (fld i) 0 vac = st i)
    (hcreneg : ∀ i, ∀ q : ℤ, q < 0 → (fld i) q vac = 0)
    (hTvac : T vac = 0)
    (hTcov : ∀ i, DCovariant T (fld i)) :
    ∃ YA : V →ₗ[ℂ] FormalDist (End ℂ V),
      (∀ i, YA (st i) = fld i) ∧
      (∀ u : V, IsField (YA u)) ∧
      (∀ u v : V, Local (YA u) (YA v)) ∧
      (∀ u : V, (YA u) 0 vac = u) ∧
      (∀ u : V, ∀ q : ℤ, q < 0 → (YA u) q vac = 0) ∧
      (∀ u : V, DCovariant T (YA u))
\end{leancode}
(\leanpath{VOA/Core/ReconstructionExistence.lean}; treated in full as Theorem~\ref{l03:thm:reconstruct}.)

\begin{dictionary}
\lean{(fld i) 0 vac} is the coefficient of $z^0$, i.e.\ the mode $(\cdot)_{(-1)}$, applied to the
vacuum; \lean{(fld i) q vac = 0} for $q<0$ says the modes $(\cdot)_{(n)}$, $n\ge0$, kill the
vacuum. So \lean{hcre0}/\lean{hcreneg} are the creation property. \lean{Submodule.span ℂ (Set.range
st) = ⊤} says the $s_i$ span $V$ (\lean{⊤} is the whole space). Note that the family
\lean{st} need not be linearly independent: the content of the theorem is precisely that if a
linear relation $\sum c_is_i=0$ holds, then $\sum c_ia^i=0$, so that $Y$ is well defined. This is
Goddard's uniqueness principle: a field local with all the $a^i$ that kills the vacuum (together
with all its $T$-derivatives) is zero.
\end{dictionary}

\begin{physics}
This is the algebraic form of the bootstrap principle: a family of mutually local fields $a^i(z)$ that are translation covariant, together with the states $s_i=a^i(0)\mathbb1$ they create, determines the whole space of local operators, namely all normal-ordered products of derivatives of the generating fields, with the OPEs fixed. Physically it is the state--operator correspondence run in reverse: from a few generating fields (the currents $\partial X$ and the vertex operators $e^{i\alpha\cdot X}$) one builds all local operators of the theory. The proof is combinatorial and does not use physical input.
\end{physics}

\begin{proof}[Proof on paper]
Define $Y$ on formal combinations by $\sum c_i\delta_i\mapsto\sum c_ia^i$. If $\sum c_is_i=0$, the
field $b=\sum c_ia^i$ is local with every $a^j$, $T$-covariant, and $b(z)\mathbb1=e^{zT}(\sum c_i
s_i)=0$ (the covariance plus $T\mathbb 1=0$ forces $b(z)\mathbb 1=e^{zT}b_{0}\mathbb1$). By
locality, $(z_1-z_2)^Nb(z_1)a^j(z_2)\mathbb 1=(z_1-z_2)^Na^j(z_2)b(z_1)\mathbb1=0$; setting
$z_2=0$ (allowed on the left since $a^j(z_2)\mathbb 1$ has no negative powers) gives
$z_1^Nb(z_1)s_j=0$, so $b$ kills every $s_j$, hence $b=0$. Thus $Y$ descends to $V$. Locality of
general $Y(u),Y(v)$ follows by bilinearity, and the remaining properties likewise.
\end{proof}

\subsection{The generating family for \texorpdfstring{$V_L$}{V\_L}}

\begin{definition}[Word states and word fields]\label{ch03:def:words}
For a list $l=[(i_1,a_1),\dots,(i_s,a_s)]$ and $\alpha\in L$, the word state is
$h_{a_1}(-i_1-1)\cdots h_{a_s}(-i_s-1)e^\alpha$ and the word field is the right-nested normal
ordered product \eqref{ch03:eq:Ygeneral}.
\end{definition}
\noindent File \leanpath{VOA/Examples/LatticeGenWordFields.lean}.
\begin{leancode}
noncomputable def wordStateLatGen : List (ℕ × Fin (rank L)) → L.carrier → LatModuleGen L
  | [], α => latChargeGen L α
  | (ia :: rest), α =>
      LinearMap.rTensor (L.carrier →₀ ℂ) (currentN (rank L) 1 ia.2 ia.1) (wordStateLatGen rest α)

noncomputable def YwordLatGenF (α : L.carrier) :
    List (ℕ × Fin (rank L)) → { f : FormalDist (End ℂ (LatModuleGen L)) // IsField f }
  | [] => ⟨Yeαvec L α, Yeαvec_isField L α⟩
  | (ia :: rest) =>
      let p := YwordLatGenF α rest
      ⟨nprod (genFieldLatGen_isField L ia.2 ia.1) p.2,
        (genFieldLatGen_isField L ia.2 ia.1).nprod p.2⟩

noncomputable def YwordLatGen (l : List (ℕ × Fin (rank L))) (α : L.carrier) :
    FormalDist (End ℂ (LatModuleGen L)) := (YwordLatGenF L α l).1
\end{leancode}

\begin{dictionary}
\lean{currentN (rank L) 1 a i} is the $z^i$-coefficient of the current, i.e.\ $h_a(-i-1)$, which
on the Fock factor is multiplication by $X_{(i,a)}$; so the word state is as stated. The subtype
\lean{\{f // IsField f\}} carries the proof that each partial product is a field, which
\lean{nprod} requires. The word states span $V_L$, because the
monomials span $\mathcal F$:
\noindent File \leanpath{VOA/Examples/LatticeGenWordFields.lean}.
\begin{leancode}
theorem wordStateLatGen_span :
    Submodule.span ℂ (Set.range (Function.uncurry (wordStateLatGen L))) = ⊤
\end{leancode}
\end{dictionary}

\begin{physics}
A word state $h_{a_1}(-i_1-1)\cdots h_{a_s}(-i_s-1)e^\alpha$ is a multi-oscillator excitation above the charged ground state $e^\alpha$: its image under the state--operator map is the normal-ordered product $:\!\partial^{i_1}J^{a_1}\cdots\partial^{i_s}J^{a_s}\,e^{i\alpha\cdot X}\!:$ of derivatives of the currents with the vertex operator, which is how every local operator of the compact boson is written. The right-nested normal ordering is the standard prescription for defining a composite operator without ambiguity. These words span $V_L$, so these fields exhaust the local operators.
\end{physics}

The translation operator must be chosen before any conformal vector is available. It is
$T=L_{-1}$, which on $V_L$ is the Fock $L_{-1}$ plus the cross term between $h_a(-1)$ and the
momentum $h_a(0)$:
\begin{leancode}
noncomputable def momGen : End ℂ (LatModuleGen L) :=
  ∑ a : Fin (rank L), (LinearMap.rTensor (L.carrier →₀ ℂ) (oscN (rank L) 1 a (-1))) * zGradeOpGen L a

noncomputable def TrecGen : End ℂ (LatModuleGen L) :=
  LinearMap.rTensor (L.carrier →₀ ℂ) (HeisenbergN.LvirN (rank L) 1 (-1)) + momGen L
\end{leancode}
(\leanpath{VOA/Examples/LatticeGenTranslation.lean}.) On paper,
$L_{-1}=\frac12\sum_a\sum_j{:}h_a(-1-j)h_a(j){:}$; the two terms with $\{j,-1-j\}=\{0,-1\}$ give
$h_a(-1)h_a(0)$, the remaining terms are the Fock $L_{-1}$ (where $h_a(0)=0$). Covariance
$[T,Y(e^\beta,z)]=\partial_zY(e^\beta,z)$ (\lean{TrecGen_Yeαvec_DCovariant}) is the identity
$\partial_zY(e^\beta,z)={:}\beta(z)Y(e^\beta,z){:}$ in disguise: the derivative of
$z^{\beta(0)}E^-E^+$ produces $\beta(0)z^{-1}$ and the creation and annihilation parts of
$\beta(z)$.

\begin{theorem}[The reconstructed vertex map]\label{ch03:thm:latYfull}
Let $L$ be even and positive-definite. There is a linear map $Y\colon V_L\to\End(V_L)[[z^{\pm1}]]$
which equals \eqref{ch03:eq:Ygeneral} on every word state and is a local, $T$-covariant field map
with the creation property.
\end{theorem}
\begin{leancode}
theorem latYfullGen_of_locality (hEven : L.Even) (hPos : L.PosDef) :
    ∃ Yfull : LatModuleGen L →ₗ[ℂ] FormalDist (End ℂ (LatModuleGen L)),
      (∀ (lα : List (ℕ × Fin (rank L)) × L.carrier),
        Yfull (Function.uncurry (wordStateLatGen L) lα)
          = Function.uncurry (fun l α => YwordLatGen L l α) lα) ∧
      (∀ u : LatModuleGen L, IsField (Yfull u)) ∧
      (∀ u v : LatModuleGen L, Local (Yfull u) (Yfull v)) ∧
      (∀ u : LatModuleGen L, (Yfull u) 0 (latVacuumGen L) = u) ∧
      (∀ u : LatModuleGen L, ∀ q : ℤ, q < 0 → (Yfull u) q (latVacuumGen L) = 0) ∧
      (∀ u : LatModuleGen L, Reconstruction.DCovariant (TrecGen L) (Yfull u))

noncomputable def YfullGen (hEven : L.Even) (hPos : L.PosDef) :
    LatModuleGen L →ₗ[ℂ] FormalDist (End ℂ (LatModuleGen L)) :=
  (latYfullGen_of_locality L hEven hPos).choose
\end{leancode}
(\leanpath{VOA/Examples/LatticeGenWordFields.lean}, \leanpath{LatticeGenVOAInstance.lean}.)

\begin{dictionary}
\lean{Function.uncurry f (l, α)} $=f\,l\,\alpha$; the first conjunct says that $Y$ of every word
state is the corresponding word field. \lean{YfullGen} is defined by \lean{.choose}, i.e.\ it is
\emph{some} map with these properties. Could a bad choice be made? No: the first conjunct fixes
$Y$ on a spanning set, and $Y$ is linear, so the map is unique. Hence \lean{YfullGen} is the
textbook vertex map \eqref{ch03:eq:Ygeneral}, not merely a map with the listed properties.
In particular $Y(e^\alpha,z)$ is \lean{Yeαvec L α} (the empty word).
\end{dictionary}

\begin{physics}
This theorem constructs the operator--state correspondence of the compactified boson: every state of $V_L$ is assigned a field $Y(u,z)$, with the generating fields being the currents $J^a(z)$ and the exponentials $:\!e^{i\alpha\cdot X}\!:$, and the assignment is local, translation covariant, and sends $u$ to $u$ on the vacuum. It does so by applying the reconstruction theorem to the oscillator currents and charged vertex operators, whose mutual locality was established in Theorem~\ref{ch03:thm:Yelocal}. It is the technical heart of the lattice construction; physically it records that the set of vertex operators of the compact boson closes under OPE.
\end{physics}

\begin{proof}[Proof on paper]
Apply Theorem~\ref{ch03:thm:reconstruct} with index set $I=\{\text{lists}\}\times L$. Mutual
locality of the word fields follows from Theorem~\ref{ch03:thm:Yelocal}, the locality of the
currents with each other and with the $Y(e^\beta)$, and Dong's lemma (normal-ordered products of
pairwise local fields are local with the same partners), by induction on the words.
$T$-covariance of the word fields follows from covariance of the generators, since
$T$-covariance is preserved by derivatives and normal-ordered products. The creation property of
a word field on the vacuum is the statement that normal-ordered products create products of
states: $({:}ab{:})(z)\mathbb1|_{z=0}=a_{(-1)}b_{(-1)}\mathbb 1$.
\end{proof}

\begin{status}
Proved, conditional only on \lean{hEven} and \lean{hPos};
\texttt{\#print axioms latYfullGen\_of\_locality} $=$ \texttt{[propext, Classical.choice,
Quot.sound]}. No use of \lean{VOA.pbw}.
\end{status}

\subsection{The conformal layer and the instance}

The conformal vector is $\omega=\frac12\sum_ah_a(-1)^2\mathbb 1$, i.e.\ the Heisenberg conformal
vector at level $1$ placed in the momentum-zero sector:
\noindent File \leanpath{VOA/Examples/LatticeGenConformal.lean}.
\begin{leancode}
noncomputable def omegaLatGen : LatModuleGen L := mkEα L 0 (HeisenbergN.ωN (rank L) 1)
\end{leancode}
If $h'_b=\sum_aQ_{ba}h_a$ with $Q^TQ=1$, then $\sum_bh'_b(-1)^2=\sum_{a,a'}(Q^TQ)_{aa'}
h_a(-1)h_{a'}(-1)=\sum_ah_a(-1)^2$; so $\omega$ does not depend on which orthonormal frame is used.
This is used in functoriality (\lean{fockMap_ωN}). The library then proves the four conformal axioms:
\noindent Files \leanpath{VOA/Examples/LatticeGenVirasoro.lean}, \leanpath{VOA/Examples/LatticeGenL0Translation.lean}, \leanpath{VOA/Examples/LatticeGenGradedField.lean}.
\begin{leancode}
theorem YfullGen_virasoro (hEven : L.Even) (hPos : L.PosDef) (m n : ℤ) :
    ⁅mode (YfullGen L hEven hPos (omegaLatGen L)) (m + 1),
      mode (YfullGen L hEven hPos (omegaLatGen L)) (n + 1)⁆
      = ((m : ℂ) - n) • mode (YfullGen L hEven hPos (omegaLatGen L)) (m + n + 1)
        + (if m + n = 0 then (((m : ℂ) ^ 3 - m) / 12 * (rank L : ℂ)) • (1 : End ℂ (LatModuleGen L))
           else 0)
theorem YfullGen_L0_grading (hEven : L.Even) (hPos : L.PosDef) {n : ℤ} (v : LatModuleGen L)
    (hv : v ∈ gradingLatGen L n) :
    mode (YfullGen L hEven hPos (omegaLatGen L)) 1 v = (n : ℂ) • v
theorem YfullGen_translation_omega (hEven : L.Even) (hPos : L.PosDef) (v : LatModuleGen L) :
    YfullGen L hEven hPos (mode (YfullGen L hEven hPos (omegaLatGen L)) 0 v)
      = deriv (YfullGen L hEven hPos v)
theorem YfullGen_grading_mode (hEven : L.Even) (hPos : L.PosDef) {p q : ℤ} (u : LatModuleGen L)
    (hu : u ∈ gradingLatGen L p) (v : LatModuleGen L) (hv : v ∈ gradingLatGen L q) (j : ℤ) :
    mode (YfullGen L hEven hPos u) j v ∈ gradingLatGen L (p + q - j - 1)
\end{leancode}
On paper: $Y(\omega,z)=\frac12\sum_a{:}h_a(z)^2{:}$ is the Sugawara field of the $r$ currents
including zero modes; on the sector $\mathcal F\otimes e^\gamma$ one has
$L_0=L_0^{\mathrm{Fock}}+\frac12\langle\gamma,\gamma\rangle$, which gives the grading, and the
Virasoro relation with $c=r$ is the Heisenberg one (Theorem~\ref{ch03:thm:virasoroN}) with the
zero modes absorbed into $(m-n)L_{m+n}$. The grading law
$\wt(u_{(j)}v)=\wt u+\wt v-j-1$ follows from the fact that $Y(e^\alpha,z)$ is homogeneous of
weight $\frac12\langle\alpha,\alpha\rangle$, via
$\frac12|\alpha+\beta|^2=\frac12|\alpha|^2+\frac12|\beta|^2+\langle\alpha,\beta\rangle$ (which uses
evenness).

\begin{theorem}[The lattice vertex operator algebra]\label{ch03:thm:latticeVOA}
For every even positive-definite lattice $L$ of rank $r$, $V_L$ with the vertex map of
Theorem~\ref{ch03:thm:latYfull}, vacuum $1\otimes e^0$, conformal vector
$\omega=\frac12\sum_ah_a(-1)^2\mathbb 1$ and the grading of Definition~\ref{ch03:def:gradingLatGen}
is a vertex operator algebra of CFT type with central charge $r$.
\end{theorem}
\begin{leancode}
noncomputable def latticeVOAGen (hEven : L.Even) (hPos : L.PosDef) :
    VertexOperatorAlgebra (LatModuleGen L) where
  grading := gradingLatGen L
  isInternal := gradingLatGen_isInternal_of_posDef L hEven hPos
  finrank := gradingLatGen_finrank_of_posDef L hEven hPos
  Y := YfullGen L hEven hPos
  vacuum := latVacuumGen L
  vacuum_mem := by
    rw [gradingLatGen_zero_eq_span_vacuum L hEven hPos]
    exact Submodule.mem_span_singleton_self (latVacuumGen L)
  grading_mode := fun u hu v hv j => YfullGen_grading_mode L hEven hPos u hu v hv j
  vacuum_Y := YfullGen_vacuum L hEven hPos
  creation := YfullGen_creation L hEven hPos
  locality := YfullGen_local L hEven hPos
  regularity := YfullGen_regularity L hEven hPos
  conformalVector := omegaLatGen L
  conformal_mem := omegaLatGen_mem L
  centralCharge := (rank L : ℂ)
  virasoro := YfullGen_virasoro L hEven hPos
  L0_grading := fun v hv => YfullGen_L0_grading L hEven hPos v hv
  translation := YfullGen_translation_omega L hEven hPos
  vacuum_grade := gradingLatGen_zero_eq_span_vacuum L hEven hPos
  positivity := fun n hn => gradingLatGen_positivity L hEven hPos n hn

theorem latticeVOAGen_centralCharge (hEven : L.Even) (hPos : L.PosDef) :
    (latticeVOAGen L hEven hPos).centralCharge = (rank L : ℂ) := rfl
\end{leancode}
(\leanpath{VOA/Examples/LatticeGenInstance.lean}.)

\begin{dictionary}
Every field of the structure is filled by a theorem stated above, and every one of those has only
\lean{hEven} and \lean{hPos} as hypotheses. The vacuum field axiom \lean{vacuum_Y} is
$Y(\mathbb 1,z)=\mathrm{id}$, which here reduces to $Y(e^0,z)=E^-(0,z)E^+(0,z)=\mathrm{id}$ (the
Schur polynomials of the zero charge vanish in positive degree). As for the Heisenberg algebra,
\lean{latticeVOAGen_centralCharge} is \lean{rfl}, but non-vacuous because the field
\lean{virasoro} is proved with exactly this value. A referee should also note what the instance
does \emph{not} contain: no Jacobi identity is listed, because in this library's axiomatics
(Theorem~\ref{l03:thm:jacobilocal}) locality plus the other axioms imply it.
\end{dictionary}

\begin{physics}
This is the Frenkel--Lepowsky--Meurman construction: the lattice VOA $V_L$ is the chiral algebra of $r$ bosons compactified on the torus $\mathbb R^r/2\pi L^*$, with $c=r$ from the free-boson Virasoro algebra and states of conformal weight $h=\frac12\langle\alpha,\alpha\rangle+$ oscillator level. Its character is $\theta_L(\tau)/\eta(\tau)^r$ with $\theta_L=\sum_{\alpha\in L}q^{\langle\alpha,\alpha\rangle/2}$, so for a unimodular $L$ it is (after the $q^{-c/24}$ shift) modular invariant up to a phase; for a general even lattice $V_L$ is a holomorphic CFT only when $L$ is unimodular. Weight-one states from norm-$2$ vectors enhance the $U(1)^r$ current algebra to a simply-laced Lie algebra (Frenkel--Kac), the stringy origin of enhanced gauge symmetry at special radii.
\end{physics}

\begin{status}
Proved, conditional only on \lean{hEven}, \lean{hPos};
\texttt{\#print axioms latticeVOAGen} $=$ \texttt{[propext, Classical.choice, Quot.sound]};
\lean{VOA.pbw} not used.
\end{status}

\section{The Leech lattice vertex operator algebra}\label{ch03:sec:leech}

\subsection{Importing the Leech lattice}

The Leech lattice $\Leech$ is constructed in the vendored library
\leanpath{VOA/LatticesAndCodes} (a formalisation of Ebeling's \emph{Lattices and Codes}; see
Chapter~\ref{ch:02}) from the extended binary Golay code, and proved there to be even,
positive-definite, unimodular, of rank $24$ and without roots. That library has its own type
\lean{_root_.IntegralLattice}; the bridge to the vertex-algebra notion is a field-by-field
repackaging:
\begin{leancode}
noncomputable def ofEbeling (Γ : _root_.IntegralLattice) : VOA.IntegralLattice.{0} where
  carrier := Γ.carrier
  form := Γ.bilin
  isSymm := ⟨fun x y => Γ.bilin_symm x y⟩

theorem ofEbeling_even (h : Γ.IsEven) : (ofEbeling Γ).Even := fun v => h v
theorem ofEbeling_posDef (h : Γ.IsPosDef) : (ofEbeling Γ).PosDef := fun v hv => h v hv
theorem ofEbeling_rank : rank (ofEbeling Γ) = Γ.rank := rfl

noncomputable abbrev leech : VOA.IntegralLattice.{0} := ofEbeling Leech.lattice

theorem leech_even : leech.Even := ofEbeling_even _ Leech.lattice_isEven
theorem leech_posDef : leech.PosDef := ofEbeling_posDef _ Leech.lattice_isPosDef
theorem leech_rank : rank leech = 24 := Leech.lattice_rank
theorem leech_no_roots (v : leech.carrier) : leech.form v v ≠ 2 := Leech.lattice_no_roots v
\end{leancode}
(\leanpath{VOA/Moonshine/LeechLatticeVOA.lean}.)

\begin{dictionary}
The proofs \lean{fun v => h v} show that evenness and positive-definiteness are \emph{the same
predicates} in both libraries, so nothing is lost or reinterpreted in the transfer; rank is even
preserved by \lean{rfl}. The four Leech facts \lean{Leech.lattice_isEven},
\lean{Leech.lattice_isPosDef}, \lean{Leech.lattice_rank}, \lean{Leech.lattice_no_roots} are
theorems of the vendored library (Chapter~\ref{ch:02} gives their statements and checks that
\lean{Leech.lattice} is the Leech lattice, i.e.\ the definition is not some other lattice that
happens to satisfy these properties; by Conway's uniqueness theorem, which is \emph{not}
formalised, these four properties together with unimodularity characterise $\Leech$ up to
isometry). The universe annotation \lean{.{0}} only fixes the size of the type.
\end{dictionary}

\begin{definition}[The Leech lattice vertex operator algebra]\label{ch03:def:leechVOA}
$V_\Leech$ is the lattice vertex operator algebra of Theorem~\ref{ch03:thm:latticeVOA} for
$L=\Leech$.
\end{definition}
\begin{leancode}
noncomputable def leechVOA : VertexOperatorAlgebra (LatGen.LatModuleGen leech) :=
  LatGen.latticeVOAGen leech leech_even leech_posDef

theorem leechVOA_centralCharge : leechVOA.centralCharge = 24 := by
  rw [leechVOA, LatGen.latticeVOAGen_centralCharge, leech_rank]; norm_num
\end{leancode}

\begin{physics}
This is the chiral algebra of the $c=24$ bosonic string compactified on the Leech torus $\mathbb R^{24}/\Lambda_{\rm Leech}$ (heuristically, the left-moving sector of the $24$-dimensional torus at the Leech point). The Leech lattice is even unimodular of rank $24$ with no vectors of norm $2$, so no momentum state has weight one and there is no enhanced gauge symmetry: the only currents are the $24$ abelian oscillator currents $\partial X^a$, and the partition function is $\theta_\Lambda/\eta^{24}=j(\tau)-744+24$ (the $24$ counting the weight-one space). The theory is a holomorphic CFT of central charge $24$ and is the starting point for the orbifold construction of the Monster module $V^\natural$.
\end{physics}

\subsection{The weight-one space}

\begin{theorem}[Rootless lattices have abelian weight-one space]\label{ch03:thm:rootless}
If $L$ is even, positive-definite and has no vectors of norm $2$, then
$(V_L)_1=\mathcal F_1\otimes e^0=\operatorname{span}\{h_a(-1)\mathbb 1\}$. In particular
$\dim(V_\Leech)_1=24$.
\end{theorem}
\begin{leancode}
theorem gradingLatGen_one_of_rootless (hEven : L.Even) (hPos : L.PosDef)
    (hR : ∀ v : L.carrier, L.form v v ≠ 2) :
    gradingLatGen L 1 = bigradα L (1, 0)

theorem leechVOA_grading_one :
    leechVOA.grading 1 = LatGen.bigradα leech (1, 0) :=
  LatGen.gradingLatGen_one_of_rootless leech leech_even leech_posDef leech_no_roots
\end{leancode}

\begin{dictionary}
\lean{bigradα L (1, 0)} is $(\mathcal F)_1\otimes e^0$, the span of the states $h_a(-1)\mathbb 1$
(Definition~\ref{ch03:def:gradingLatGen}). The hypothesis \lean{hR} is ``no roots''. The
statement is an equality of submodules, so it says both inclusions; it is the precise form of the
statement ``$V_\Leech$ has no weight-one currents other than the Heisenberg ones''. The dimension
$24$ is then the Fock count $\dim\mathcal F_1=r$.
\end{dictionary}

\begin{physics}
For a rootless lattice, the weight-one states are only oscillators, because the exponentials $:\!e^{i\alpha\cdot X}\!:$ have weight $\langle\alpha,\alpha\rangle/2\ge2$, the minimal nonzero norm being at least $4$. So the current algebra is the abelian $U(1)^{24}$: the Leech torus compactification has no non-abelian gauge symmetry, only the $24$ translations of the torus, and the dimension $24$ equals the number of independent boson currents. In Niemeier language, the Leech lattice is the unique one of the $24$ unimodular lattices with no roots, hence with the trivial (abelian) current algebra.
\end{physics}

\begin{proof}[Proof on paper]
A summand $(\mathcal F)_{1-|\alpha|^2/2}\otimes e^\alpha$ with $\alpha\neq0$ has
$|\alpha|^2/2\geq1$ by positivity and evenness, and $|\alpha|^2/2\neq1$ because $|\alpha|^2\ne2$
(no roots), so
$|\alpha|^2/2\geq2$ and the Fock weight is negative: the summand is zero. Only $\alpha=0$ remains.
\end{proof}

\begin{status}
Proved, unconditional (all hypotheses discharged for $\Leech$);
\texttt{\#print axioms} of \lean{leechVOA}, \lean{leechVOA_centralCharge},
\lean{leechVOA_grading_one} $=$ \texttt{[propext, Classical.choice, Quot.sound]}. The dimensions
$\dim(V_\Leech)_2=196884$ and $\dim(V_\Leech)_3=21493760$ are proved elsewhere
(\lean{leechVOA_finrank_grading_two}, \lean{leechVOA_finrank_grading_three}; Chapter~\ref{ch:05}).
\end{status}

\begin{remark}[The physical meaning of $V_\Leech$]
$V_\Leech$ is the chiral algebra of $24$ free bosons compactified on $\chthreeR^{24}/\Leech$. Because
$\Leech$ is unimodular, $V_\Leech$ is \emph{holomorphic}, in the following sense. A vertex
operator algebra $V$ is called holomorphic if every (admissible) $V$-module is a direct sum of
copies of the adjoint module $V$; for a \emph{rational} $V$ (one whose admissible modules are
completely reducible, with finitely many irreducibles) this is the same as saying that $V$ is its
only irreducible module. For lattice vertex operator algebras this is governed by the
\emph{dual lattice} $L^*=\{\lambda\in L\otimes\Q:\langle\lambda,L\rangle\subseteq\Z\}$: by a theorem
of Dong (C.~Dong, ``Vertex algebras associated with even lattices'', \emph{J.\ Algebra}~161
(1993); see also Lepowsky--Li, Chapter~6), $V_L$ is rational and its irreducible modules are
exactly the spaces $V_{L+\lambda}=\mathcal F\otimes\C[L+\lambda]$, one for each coset
$\lambda+L\in L^*/L$. The lattice is unimodular precisely when $L^*=L$ (Chapter~\ref{ch:02}), so then
$V_L$ itself is the only irreducible module and $V_L$ is holomorphic. \emph{Neither Dong's
classification nor the holomorphy of $V_\Leech$ is proved in the library}; nothing in the
Moonshine construction depends on them, and they are mentioned here only to explain the picture.
Because $\Leech$ has no roots, the $24$ currents do not enhance to a non-abelian affine algebra. Its
character is $q^{-1}\theta_\Leech(q)/\eta(q)^{24}=J(q)+24$, so it differs from the Moonshine
module only by the $24$ states $h_a(-1)\mathbb 1$; removing them is the purpose of the
$\Z/2$-orbifold of Chapters~\ref{ch:04}--\ref{ch:10}.
\end{remark}

\section{Sign automorphisms \texorpdfstring{$g_\chi$}{g\_chi}}\label{ch03:sec:sign}

\subsection{The involution \texorpdfstring{$\theta$}{theta}, recalled}

The automorphism $\theta$ of $V_L$ lifting $-1\in O(L)$ is the subject of Chapter~\ref{ch:04}; we
need its definition because the theorems below mention it. Let $\sigma$ be the algebra
automorphism of the Fock space that negates every variable, $\sigma(X_{(i,a)})=-X_{(i,a)}$; on a
monomial of degree $d$ it acts by $(-1)^d$. The reason: $-1\in O(L)$ acts on
$\mathfrak h=L\otimes\C$ by $h\mapsto-h$, hence on the Heisenberg algebra by $h(n)\mapsto-h(n)$,
and the induced action on the Fock space sends $h_{a_1}(-n_1)\cdots h_{a_d}(-n_d)\mathbb1$ to
$(-1)^d$ times itself, i.e.\ negates every variable. Then
$\theta(p\otimes e^\alpha)=\sigma(p)\otimes e^{-\alpha}$.
\begin{leancode}
noncomputable def sigmaAlg : FockN r →ₐ[ℂ] FockN r := MvPolynomial.aeval (fun v => -(X v : FockN r))

noncomputable def sigmaN : FockN r ≃ₗ[ℂ] FockN r :=
  LinearEquiv.ofInvolutive (sigmaAlg r).toLinearMap (sigmaAlg_involutive r)

noncomputable def thetaLat : LatModuleGen L ≃ₗ[ℂ] LatModuleGen L :=
  TensorProduct.congr (sigmaN (rank L)) (Finsupp.domLCongr (Equiv.neg L.carrier))

theorem thetaLat_latStateGen (p : FockN (rank L)) (α : L.carrier) :
    thetaLat L (latStateGen L p α) = latStateGen L (sigmaN (rank L) p) (-α)
\end{leancode}
(\leanpath{VOA/Moonshine/LatticeInvolution.lean}.) \lean{aeval f} is the algebra map determined
by sending each variable $v$ to $f(v)$; \lean{Finsupp.domLCongr (Equiv.neg _)} relabels
$e^\alpha\mapsto e^{-\alpha}$. Chapter~\ref{ch:04} proves that $\theta$ is an automorphism,
\lean{thetaAut : Aut (latticeVOAGen L hEven hPos)}; this uses $\varepsilon(-\alpha,-\beta)=
\varepsilon(\alpha,\beta)$ (Remark~\ref{ch03:rem:negneg}).

\subsection{The sign operator}

\begin{definition}[Sign automorphism]\label{ch03:def:chiLat}
Let $\chi\colon L\to\Z/2$ be a group homomorphism. Define $g_\chi(p\otimes e^\alpha)=
(-1)^{\chi(\alpha)}p\otimes e^\alpha$.
\end{definition}
\begin{leancode}
noncomputable def sgn (z : ZMod 2) : ℂ := (-1 : ℂ) ^ z.val

noncomputable def chiSign (α : L.carrier) : ℂ := VOA.Moonshine.Extraspecial.sgn (χ α)

noncomputable def chiGroupMap : (L.carrier →₀ ℂ) →ₗ[ℂ] (L.carrier →₀ ℂ) :=
  Finsupp.linearCombination ℂ (fun α => chiSign L χ α • eGen L α)

noncomputable def chiGroupEquiv : (L.carrier →₀ ℂ) ≃ₗ[ℂ] (L.carrier →₀ ℂ) :=
  LinearEquiv.ofInvolutive (chiGroupMap L χ) (chiGroupMap_involutive L χ)

noncomputable def chiLat : LatModuleGen L ≃ₗ[ℂ] LatModuleGen L :=
  TensorProduct.congr (LinearEquiv.refl ℂ (FockN (rank L))) (chiGroupEquiv L χ)

theorem chiLat_latStateGen (p : FockN (rank L)) (α : L.carrier) :
    chiLat L χ (latStateGen L p α) = chiSign L χ α • latStateGen L p α
\end{leancode}
(\leanpath{VOA/Moonshine/FrameSign.lean}, with section variable \lean{(χ : L.carrier →+ ZMod 2)};
\lean{sgn} is in \leanpath{VOA/Moonshine/ExtraspecialModule.lean}.)

\begin{dictionary}
\lean{L.carrier →+ ZMod 2} is the type of additive group homomorphisms $L\to\Z/2$.
\lean{z.val} is the representative of $z\in\Z/2$ in $\{0,1\}$, so \lean{sgn z} $=(-1)^z$.
\lean{LinearEquiv.ofInvolutive} turns a linear map $f$ with $f\circ f=\mathrm{id}$ into a linear
equivalence. \lean{TensorProduct.congr f g} is $f\otimes g$. So \lean{chiLat} $=\mathrm{id}\otimes
g_\chi$, matching the definition; \lean{chiLat_latStateGen} is its defining formula.
\end{dictionary}

\begin{physics}
The sign automorphism $g_\chi$ multiplies a momentum-$\alpha$ state by $(-1)^{\chi(\alpha)}$. It is the element $e^{i\pi\,\chi\cdot\hat p}$ of the target-space translation group $U(1)^r$ (a half-period shift along the torus in the direction dual to $\chi$), and its action on the charged sectors is the standard $\Z/2$ phase symmetry. Sign automorphisms of this type are the ones that appear as the lifts of the frame signs in the construction of the Monster: they distinguish the momentum sectors by a parity.
\end{physics}

\begin{theorem}[$g_\chi$ is an automorphism of order $\le2$ commuting with $\theta$]\label{ch03:thm:chiAut}
Let $L$ be even and positive-definite and $\chi\colon L\to\Z/2$ a homomorphism. Then $g_\chi\in
\Aut(V_L)$, $g_\chi^2=1$ and $g_\chi\theta=\theta g_\chi$.
\end{theorem}
\begin{leancode}
noncomputable def chiAut : Aut (latticeVOAGen L hEven hPos) where
  toLinearEquiv := chiLat L χ
  map_smul' u n v := chiLat_mode L χ hEven hPos u n v
  map_vacuum := chiLat_vacuum L χ
  map_omega := chiLat_omega L χ

theorem chiAut_sq : chiAut L χ hEven hPos * chiAut L χ hEven hPos = 1

theorem chiLat_thetaLat (x : LatModuleGen L) :
    chiLat L χ (thetaLat L x) = thetaLat L (chiLat L χ x)

theorem chiAut_mul_thetaAut :
    chiAut L χ hEven hPos * thetaAut L hEven hPos = thetaAut L hEven hPos * chiAut L χ hEven hPos
\end{leancode}

\begin{dictionary}
\lean{chiLat_mode} is the statement $g_\chi(u_{(n)}v)=(g_\chi u)_{(n)}(g_\chi v)$ for all
$u,n,v$, the content of \lean{map_smul'}. The product \lean{*} in \lean{Aut} is composition and
\lean{1} the identity. Note that $g_\chi=\mathrm{id}$ is possible (for $\chi=0$); the theorem does
not claim $g_\chi\ne1$, and nothing downstream needs it.
\end{dictionary}

\begin{physics}
This is the statement that a half-shift on the torus is a symmetry of the theory: the $U(1)^r$ momentum phase $(-1)^{\chi(\alpha)}$ is a character of the momentum lattice, hence is compatible with the OPE since momentum is conserved ($\chi(\alpha+\beta)=\chi(\alpha)+\chi(\beta)$) and gives a $\Z/2$ symmetry of the VOA. It commutes with the lattice involution $\theta$ ($\alpha\mapsto-\alpha$, $X\mapsto-X$, the T-duality or reflection symmetry of the torus) because the sign is invariant under $\alpha\mapsto-\alpha$, as $\chi$ takes values in $\Z/2$. Such automorphisms are the orbifold data used to build new holomorphic CFTs.
\end{physics}

\begin{proof}[Proof on paper]
By the uniqueness built into the reconstruction (Theorem~\ref{ch03:thm:latYfull}), it suffices to
check $g_\chi Y(w,z)g_\chi^{-1}=Y(g_\chi w,z)$ on generators. The currents $h_a(z)$, including the
zero mode, preserve each momentum sector, so commute with $g_\chi$. The field $Y(e^\alpha,z)$
maps the $\beta$-sector to the $(\alpha+\beta)$-sector, so
$g_\chi Y(e^\alpha,z)=(-1)^{\chi(\alpha+\beta)-\chi(\beta)}Y(e^\alpha,z)g_\chi=(-1)^{\chi(\alpha)}
Y(e^\alpha,z)g_\chi$ by additivity of $\chi$. Normal-ordered products multiply the scalars, so
$g_\chi Y_{\mathrm{word}}(l,\alpha)g_\chi^{-1}=(-1)^{\chi(\alpha)}Y_{\mathrm{word}}(l,\alpha)$,
which matches $g_\chi(\text{word state})=(-1)^{\chi(\alpha)}\text{word state}$. By linearity over
the spanning word states, $g_\chi$ intertwines all vertex operators. It fixes $\mathbb 1$ and
$\omega$, which have momentum $0$. Commuting with $\theta$: $\chi(-\alpha)=\chi(\alpha)$ in $\Z/2$.
\end{proof}

\begin{theorem}[Eigenspaces of $g_\chi$ by momentum]\label{ch03:thm:chieig}
For $x\in V_L$: $g_\chi x=x$ iff every momentum component $x_\gamma$ with $\chi(\gamma)\ne0$
vanishes; $g_\chi x=-x$ iff every $x_\gamma$ with $\chi(\gamma)=0$ vanishes.
\end{theorem}
\noindent File \leanpath{VOA/Moonshine/FrameSign.lean}.
\begin{leancode}
theorem chiLat_eq_self_iff (x : LatModuleGen L) :
    chiLat L χ x = x ↔ ∀ α, χ α ≠ 0 → coordEα L α x = 0

theorem chiLat_eq_neg_iff (x : LatModuleGen L) :
    chiLat L χ x = -x ↔ ∀ α, χ α = 0 → coordEα L α x = 0
\end{leancode}
\begin{dictionary}
\lean{coordEα L α x} is the momentum-$\alpha$ component $x_\alpha\in\mathcal F$ of $x$
(Section~\ref{ch03:sec:latticespace}). Thus the fixed points of $g_\chi$ form exactly the
subspace $\mathcal F\otimes\C[\ker\chi]$ (as a vector space it looks like $V_{\ker\chi}$, but no
claim about the vertex algebra $V_{\ker\chi}$ is made here), and the $-1$ eigenspace is
$\mathcal F\otimes\C[L\setminus\ker\chi]$. These are pure linear algebra (no hypotheses on $L$).
\end{dictionary}

\begin{physics}
The $\pm1$ eigenspaces of $g_\chi$ are the sectors of momentum with $\chi(\gamma)=0$ and $\chi(\gamma)=1$: the momentum lattice splits into the sublattice $\ker\chi$ of index $2$ and its nontrivial coset, and $V_L=V_{\ker\chi}\oplus V_{\text{coset}}$ accordingly. In an orbifold by $g_\chi$ these are the untwisted states kept (the even sector) and removed (the odd sector); the invariant subalgebra $V_L^{g_\chi}$ is the lattice algebra $V_{\ker\chi}$ of the index-two sublattice, plus twisted sectors.
\end{physics}
\begin{proof}[Proof on paper]
Compare components: $(g_\chi x)_\alpha=(-1)^{\chi(\alpha)}x_\alpha$, and $x_\alpha=-x_\alpha$
forces $x_\alpha=0$ in characteristic $0$. A vector is determined by its components.
\end{proof}
\begin{status}
\lean{chiAut}, \lean{chiAut_sq}, \lean{chiAut_mul_thetaAut}: proved, conditional only on
\lean{hEven}, \lean{hPos}. \lean{chiLat_eq_self_iff}, \lean{chiLat_eq_neg_iff},
\lean{chiLat_thetaLat}: proved, no hypotheses. \texttt{\#print axioms} (re-checked):
\texttt{[propext, Classical.choice, Quot.sound]}.
\end{status}

In the Moonshine construction (Chapters~\ref{ch:11}--\ref{ch:13}), $g_\chi$ for $\chi$ the
character with kernel $\Leech_0$ (an index-$2$ sublattice) splits $V_\Leech^+$ into the pieces
$V_0^+\oplus V_1^+$ used by the frame/triality argument; the fact that this splitting is by
automorphism eigenspaces makes the products $u_{(n)}v$ respect it for free.

\section{Functoriality of lattice vertex operator algebras}\label{ch03:sec:functoriality}

\subsection{The problem}

An isometric embedding $\iota\colon K\to L$ of even lattices should induce a homomorphism of
vertex operator algebras $V_K\to V_L$ sending $e^\alpha\mapsto e^{\iota\alpha}$ and
$\alpha(-n)\mapsto(\iota\alpha)(-n)$. In the library this is not automatic, for two reasons, both
coming from the choice of bases:
\begin{enumerate}
\item the Fock spaces of $V_K$ and $V_L$ are written in orthonormal oscillators relative to
  $\sqrt{G_K}$ and $\sqrt{G_L}$ (Section~\ref{ch03:sec:frame}); the map on oscillators is an
  orthogonal matrix $Q$ that must be computed;
\item the cocycles $\varepsilon_K$ and $\varepsilon_L\circ(\iota\times\iota)$ are built from
  unrelated bases, so they differ; one must correct $e^\alpha\mapsto c(\alpha)e^{\iota\alpha}$ by a
  function $c$ with $\varepsilon_L(\iota\alpha,\iota\beta)c(\alpha)c(\beta)=\varepsilon_K(\alpha,
  \beta)c(\alpha+\beta)$.
\end{enumerate}
The files \leanpath{VOA/Moonshine/LatticeMorphFock.lean}, \leanpath{LatticeMorph.lean} and
\leanpath{LatticeMorphRigid.lean} solve both, and the solution of the second contains a genuine
design point: the correction $c$ must be allowed to take values in the fourth roots of unity.

\subsection{Orthogonal change of oscillator variables}

\begin{definition}\label{ch03:def:fockMap}
For a complex $s\times r$ matrix $Q$, let $\varphi_Q\colon\mathcal F_r\to\mathcal F_s$ be the
algebra homomorphism with $\varphi_Q(X_{(k,a)})=\sum_bQ_{ba}X_{(k,b)}$ (the same matrix at every
mode level $k$).
\end{definition}
\begin{leancode}
noncomputable def fockMap : FockN r →ₐ[ℂ] FockN s :=
  MvPolynomial.aeval (fun v : ℕ × Fin r => ∑ b : Fin s, Q b v.2 • (X (v.1, b) : FockN s))

theorem fockMap_oscN (hQ : Qᵀ * Q = 1) (k : ℂ) (a : Fin r) (n : ℤ) (p : FockN r) :
    fockMap Q (oscN r k a n p) = ∑ b : Fin s, Q b a • oscN s k b n (fockMap Q p)

theorem fockMap_ωN (hQ' : Q * Qᵀ = 1) : fockMap Q (ωN r 1) = ωN s 1
\end{leancode}
\begin{dictionary}
\lean{Qᵀ * Q = 1} says the columns of $Q$ are orthonormal for the \emph{bilinear} (not Hermitian)
dot product; this is what preserves $[h_a(m),h_b(n)]=m\delta_{ab}\delta_{m+n,0}$. It is needed only
for annihilators (chain rule: $\partial_{X_{(k,b)}}\varphi_Q=\sum_aQ_{ba}\varphi_Q\partial_{X_{(k,a)}}$
requires $Q^TQ=1$ to invert). \lean{Q * Qᵀ = 1} (rows orthonormal) is needed for
$\varphi_Q(\omega)=\omega$, since $\sum_a(\sum_bQ_{ba}X_b)^2=\sum_{b,b'}(QQ^T)_{bb'}X_bX_{b'}$.
When $s>r$, only the first holds; this is why a smaller-rank embedding does not send $\omega$ to
$\omega$ (we say such a homomorphism is not \emph{conformal}).
\end{dictionary}

\begin{physics}
An isometry $\iota\colon K\to L$ of lattices, written in orthonormal coordinates, is an $r\times s$ real matrix $Q$ acting on the boson fields: it is a rotation or embedding of the target space (a linear change of the fields $X^a\mapsto Q^a_{\ b}X^b$). The map $\varphi_Q$ is the induced action on the oscillator Fock space, applying the same linear map to every mode $\alpha_{-n}$ because the current is a single vector-valued field. This is how a symmetry or inclusion of target tori acts on the free-field oscillators.
\end{physics}

The matrix attached to an isometry is $Q=\sqrt{G_L}\,M_\iota\,\sqrt{G_K}^{-1}$, where $M_\iota$ is
the integer matrix of $\iota$ in the chosen bases:
\begin{leancode}
noncomputable def latMatrix (ι : K.carrier →ₗ[ℤ] L.carrier) :
    Matrix (Fin (latRank L)) (Fin (latRank K)) ℝ :=
  Matrix.of fun j i => coordR L (ι (latticeBasis K i)) j

noncomputable def frameMatrix (ι : K.carrier →ₗ[ℤ] L.carrier) :
    Matrix (Fin (latRank L)) (Fin (latRank K)) ℝ :=
  frame L * latMatrix K L ι * (frame K)⁻¹

noncomputable def oscMatrix (ι : K.carrier →ₗ[ℤ] L.carrier) :
    Matrix (Fin (rank L)) (Fin (rank K)) ℂ :=
  ((frameMatrix K L ι).map Complex.ofRealHom).submatrix (rankEquiv L) (rankEquiv K)

theorem oscMatrix_chargeCoord (hK : K.PosDef) (ι : K.carrier →ₗ[ℤ] L.carrier) (α : K.carrier) :
    oscMatrix K L ι *ᵥ chargeCoord K α = chargeCoord L (ι α)

theorem oscMatrix_transpose_mul (hK : K.PosDef) (hL : L.PosDef)
    (ι : K.carrier →ₗ[ℤ] L.carrier) (hι : ∀ α β, L.form (ι α) (ι β) = K.form α β) :
    (oscMatrix K L ι)ᵀ * oscMatrix K L ι = 1

theorem oscMatrix_mul_transpose (hK : K.PosDef) (hL : L.PosDef)
    (ι : K.carrier →ₗ[ℤ] L.carrier) (hι : ∀ α β, L.form (ι α) (ι β) = K.form α β)
    (hr : rank K = rank L) :
    oscMatrix K L ι * (oscMatrix K L ι)ᵀ = 1
\end{leancode}
\begin{dictionary}
\lean{K.carrier →ₗ[ℤ] L.carrier} is a $\Z$-linear map; the hypothesis \lean{hι} says it is an
isometry, $\langle\iota\alpha,\iota\beta\rangle_L=\langle\alpha,\beta\rangle_K$ (isometries of
positive-definite lattices are automatically injective, \lean{isometry_injective}).
\lean{oscMatrix_chargeCoord} says $Q\alpha^\perp=(\iota\alpha)^\perp$, i.e.\ $\varphi_Q$ sends
$\alpha(-n)$ to $(\iota\alpha)(-n)$. $Q^TQ=\sqrt{G_K}^{-1}M_\iota^TG_LM_\iota\sqrt{G_K}^{-1}=
\sqrt{G_K}^{-1}G_K\sqrt{G_K}^{-1}=1$ because $M_\iota^TG_LM_\iota=G_K$ is the isometry property in
coordinates; with equal ranks a one-sided inverse is two-sided.
\end{dictionary}
(A module docstring of \leanpath{LatticeMorphFock.lean} mentions a declaration
\lean{LatGen.exists_fockMatrix}; no declaration of that name exists. Its content is realised by
\lean{oscMatrix} and the three lemmas above.)

\subsection{Comparing two cocycles: the coboundary lemma}

\begin{lemma}[Coboundary lemma]\label{ch03:lem:cobSign}
Let $M$ be a free $\Z$-module with basis $(B_i)_{i<n}$ and $D\colon M\times M\to\Z$ bilinear with
$D(x,y)+D(y,x)$ even for all $x,y$. Let $S$ be the symmetric bilinear form with matrix
$S_{ij}=D(B_i,B_j)$ for $i\le j$ and $S_{ij}=D(B_j,B_i)$ for $i>j$, and let
$c(\alpha)=i^{S(\alpha,\alpha)}$. Then $c(0)=1$, $c(-\alpha)=c(\alpha)$, $c(\alpha)\neq0$ and
\begin{equation}\label{ch03:eq:cob}
  c(\alpha+\beta)=c(\alpha)\,c(\beta)\,(-1)^{D(\alpha,\beta)} .
\end{equation}
\end{lemma}
\begin{leancode}
noncomputable def symMatrix : Matrix (Fin n) (Fin n) ℤ :=
  fun i j => if i ≤ j then D (B i) (B j) else D (B j) (B i)

noncomputable def symForm : M →ₗ[ℤ] M →ₗ[ℤ] ℤ := Matrix.toLinearMap₂ B B (symMatrix B D)

noncomputable def cobSign (α : M) : ℂ := Complex.I ^ (symForm B D α α)

theorem cobSign_zero : cobSign B D 0 = 1
theorem cobSign_neg (α : M) : cobSign B D (-α) = cobSign B D α
theorem cobSign_add (hD : ∀ x y, 2 ∣ D x y + D y x) (α β : M) :
    cobSign B D (α + β) = cobSign B D α * cobSign B D β * (-1 : ℂ) ^ (D α β)
\end{leancode}
\begin{dictionary}
\lean{Complex.I} is $i=\sqrt{-1}$ and \lean{Complex.I ^ (k : ℤ)} its integer power, so $c$ takes
values in $\{1,i,-1,-i\}$. \lean{Matrix.toLinearMap₂ B B S} is the bilinear form with matrix $S$
in the basis $B$. (The module docstring of \leanpath{LatticeMorphFock.lean} describes $c$ as
``$\pm1$-valued'' with exponent $\sum_iD_{ii}a_i(a_i-1)+2\sum_{i<j}D_{ij}a_ia_j$; that description is
stale. The code, which is what the theorems are about, uses $S(\alpha,\alpha)=\sum_iD_{ii}a_i^2+
2\sum_{i<j}D_{ij}a_ia_j$ and is genuinely $\Z/4$-valued; see Remark~\ref{ch03:rem:Z4}.)
\end{dictionary}

\begin{physics}
This lemma is a purely arithmetic statement: the quadratic function $c(\alpha)=i^{S(\alpha,\alpha)}$ is a coboundary that converts a bilinear $\Z/2$-cocycle with symmetric part $D+D^T$ into a trivial one, so that two cocycle choices differ by a function $c$. In the physics of Klein factors, it says that the ambiguity in the sign conventions $\varepsilon$ can be absorbed by a redefinition of the momentum states $e^\alpha\mapsto c(\alpha)e^\alpha$, i.e.\ it is a gauge transformation. It is the technical ingredient in the next theorem.
\end{physics}
\begin{proof}[Proof on paper]
$S(\alpha+\beta,\alpha+\beta)=S(\alpha,\alpha)+S(\beta,\beta)+2S(\alpha,\beta)$ by symmetry, and
$i^{2S(\alpha,\beta)}=(-1)^{S(\alpha,\beta)}$. It remains to see $S(\alpha,\beta)\equiv D(\alpha,
\beta)\pmod2$: the matrices differ only below the diagonal, by $D(B_j,B_i)-D(B_i,B_j)\equiv
D(B_j,B_i)+D(B_i,B_j)\equiv0$. Finally $S(-\alpha,-\alpha)=S(\alpha,\alpha)$ gives $c(-\alpha)=
c(\alpha)$.
\end{proof}

Apply this to the discrepancy between the two cocycle exponents:
\begin{leancode}
noncomputable def cocForm : L.carrier →ₗ[ℤ] L.carrier →ₗ[ℤ] ℤ := ...   -- cocExp as a bilinear map

noncomputable def cocDiff (ι : K.carrier →ₗ[ℤ] L.carrier) : K.carrier →ₗ[ℤ] K.carrier →ₗ[ℤ] ℤ :=
  (cocForm L).compl₁₂ ι ι - cocForm K

theorem cocDiff_symm_even (hK : K.Even) (hL : L.Even) (ι : K.carrier →ₗ[ℤ] L.carrier)
    (hι : ∀ α β, L.form (ι α) (ι β) = K.form α β) (x y : K.carrier) :
    2 ∣ cocDiff K L ι x y + cocDiff K L ι y x

noncomputable def latCob (ι : K.carrier →ₗ[ℤ] L.carrier) (α : K.carrier) : ℂ :=
  cobSign (latticeBasis K) (cocDiff K L ι) α
\end{leancode}
(\lean{cocForm} packages \lean{cocExp L (latticeBasis L)} as a bilinear map; we have elided the
bilinearity proofs, and \lean{latticeCocycle_eq_zpow} records
$\varepsilon_L(\alpha,\beta)=(-1)^{\lean{cocForm}(\alpha,\beta)}$.)

\begin{theorem}[Cocycle comparison]\label{ch03:thm:cocCompare}
If $K,L$ are even and $\iota$ is an isometry, then $c=\lean{latCob}$ satisfies
$\varepsilon_L(\iota\alpha,\iota\beta)\,c(\alpha)c(\beta)=\varepsilon_K(\alpha,\beta)\,c(\alpha+
\beta)$, and $c(0)=1$, $c\ne0$, $c(-\alpha)=c(\alpha)$.
\end{theorem}
\begin{leancode}
theorem latticeCocycle_map_mul_latCob (hK : K.Even) (hL : L.Even) (ι : K.carrier →ₗ[ℤ] L.carrier)
    (hι : ∀ α β, L.form (ι α) (ι β) = K.form α β) (α β : K.carrier) :
    latticeCocycle L (ι α) (ι β) * (latCob K L ι α * latCob K L ι β)
      = latticeCocycle K α β * latCob K L ι (α + β)
theorem latCob_ne_zero (ι : K.carrier →ₗ[ℤ] L.carrier) (α : K.carrier) : latCob K L ι α ≠ 0
theorem latCob_zero (ι : K.carrier →ₗ[ℤ] L.carrier) : latCob K L ι 0 = 1
theorem latCob_neg (ι : K.carrier →ₗ[ℤ] L.carrier) (α : K.carrier) :
    latCob K L ι (-α) = latCob K L ι α
\end{leancode}

\begin{physics}
This says that two different Klein-factor gauges, $\varepsilon_K$ on the source lattice and $\varepsilon_L\circ(\iota\times\iota)$ on the target, differ by a coboundary $c$: a relabelling $e^\alpha\mapsto c(\alpha)e^\alpha$ of the momentum eigenstates. Since only cohomology classes of the cocycle are physically meaningful, this is the statement that an isometric embedding of lattices respects the Klein-factor structure up to the gauge freedom. The correction $c$ is the phase needed so that the vertex operators $:\!e^{i\alpha\cdot X}\!:$ of $V_K$ are carried to those of $V_L$.
\end{physics}
\begin{proof}[Proof on paper]
$D=\chthreecoc_L\circ(\iota\times\iota)-\chthreecoc_K$ is bilinear, and
$D(x,y)+D(y,x)\equiv\langle\iota x,\iota y\rangle-\langle x,y\rangle=0\pmod 2$ by
Theorem~\ref{ch03:thm:cocycle}(iv)--(v) for both lattices and the isometry property. Lemma
\ref{ch03:lem:cobSign} gives $c(\alpha+\beta)=c(\alpha)c(\beta)(-1)^{D(\alpha,\beta)}$, i.e.\
$c(\alpha+\beta)\varepsilon_K(\alpha,\beta)=c(\alpha)c(\beta)\varepsilon_L(\iota\alpha,\iota\beta)$
using $(-1)^{-x}=(-1)^x$.
\end{proof}

\begin{remark}[Why $\Z/4$ and not $\pm1$: the design point]\label{ch03:rem:Z4}
Any function $c$ satisfying \eqref{ch03:eq:cob} obeys, at $\beta=-\alpha$,
$1=c(0)=c(\alpha)c(-\alpha)(-1)^{D(\alpha,\alpha)}$. If $c$ were $\pm1$-valued this would give
\begin{equation*}
  c(-\alpha)=c(\alpha)\,(-1)^{D(\alpha,\alpha)} .
\end{equation*}
The function $\alpha\mapsto(-1)^{D(\alpha,\alpha)}$ is a character of $K$ (because
$D(\alpha,\beta)+D(\beta,\alpha)$ is even), and it is in general \emph{nontrivial}. The morphism
$\Phi(p\otimes e^\alpha)=c(\alpha)\varphi(p)\otimes e^{\iota\alpha}$ then satisfies
$\Phi\theta_K(p\otimes e^\alpha)=c(-\alpha)\,\sigma\varphi(p)\otimes e^{-\iota\alpha}$ and
$\theta_L\Phi(p\otimes e^\alpha)=c(\alpha)\,\sigma\varphi(p)\otimes e^{-\iota\alpha}$, so
$\Phi\theta_K=\theta_L\Phi\,g_\chi$, where $g_\chi$ is the sign automorphism of $V_K$
(Section~\ref{ch03:sec:sign}) for the character $\chi(\alpha)=D(\alpha,\alpha)\bmod 2$: $\Phi$
commutes with $\theta$ \emph{only up to a character}. The orbifold argument needs
$\Phi\theta_K=\theta_L\Phi$ on the nose, because it transports the fixed-point subalgebra
$V_K^+$. Allowing $c(\alpha)\in\{\pm1,\pm i\}$ solves $c(\alpha)^2=(-1)^{D(\alpha,\alpha)}$ with
$c(-\alpha)=c(\alpha)$, which is exactly what $i^{S(\alpha,\alpha)}$ does.

\emph{A concrete example.} Let $K=L=A_2$ with basis $B_0,B_1$ as in Example~\ref{ch03:ex:A1A2},
so $\chthreecoc(\alpha,\beta)=-\alpha_1\beta_0$, and let $\iota$ be the reflection in $B_1$:
$\iota B_0=B_0+B_1$, $\iota B_1=-B_1$. Then $D(B_0,B_0)=\chthreecoc(B_0+B_1,B_0+B_1)-0=-1$ is odd.
So every $\pm1$-valued solution $c'$ has $c'(-B_0)=-c'(B_0)$, and the $\pm1$ version of $\Phi$
fails to commute with $\theta$. The library's $c$ has $S_{00}=D_{00}=-1$, so $c(B_0)=i^{-1}=-i=
c(-B_0)$. (Here we have used the basis order shown; for the library's \lean{chooseBasis} the order
is not known, but the phenomenon is basis-independent: whenever $D(\alpha,\alpha)$ is odd for some
$\alpha$, a $\pm1$-valued coboundary cannot be even.)
\end{remark}

\begin{warning}
Recorded in \leanpath{Notes/M8GluingPlan.md} (Status 2026-10-06, ``after the functoriality
bridge''): the cocycle comparison \emph{must} use the $\Z/4$-valued coboundary; a $\pm1$-valued
one would commute with $\theta$ only up to a character. A future refactor that ``simplifies''
\lean{cobSign} to a sign would silently break \lean{latMorph_thetaLat} and everything downstream
that identifies $V_{\Leech_0}^+$ with a subalgebra of $V_\Leech^+$.
\end{warning}

\subsection{The morphism \texorpdfstring{$\Phi_\iota$}{Phi\_iota}}

\begin{definition}\label{ch03:def:latMorph}
For a function $c\colon K\to\C$ let $\Phi^c(p\otimes e^\alpha)=c(\alpha)\,\varphi_Q(p)\otimes
e^{\iota\alpha}$ with $Q=\lean{oscMatrix K L ι}$, and $\Phi_\iota:=\Phi^{\lean{latCob}}$.
\end{definition}
\noindent File \leanpath{VOA/Moonshine/LatticeMorph.lean}.
\begin{leancode}
noncomputable def momMap (c : K.carrier → ℂ) : (K.carrier →₀ ℂ) →ₗ[ℂ] (L.carrier →₀ ℂ) :=
  Finsupp.linearCombination ℂ (fun α => c α • eGen L (ι α))

noncomputable def latMorphC (c : K.carrier → ℂ) : LatModuleGen K →ₗ[ℂ] LatModuleGen L :=
  TensorProduct.map (fockMap (oscMatrix K L ι)).toLinearMap (momMap K L ι c)

noncomputable def latMorph : LatModuleGen K →ₗ[ℂ] LatModuleGen L := latMorphC K L ι (latCob K L ι)

theorem latMorph_latStateGen (p : FockN (rank K)) (α : K.carrier) :
    latMorph K L ι (latStateGen K p α)
      = latCob K L ι α • latStateGen L (fockMap (oscMatrix K L ι) p) (ι α)
\end{leancode}

\begin{physics}
The map $\Phi_\iota$ is the induced action on the chiral algebra of an isometric embedding $\iota\colon K\hookrightarrow L$ of the momentum lattices: it rotates the oscillators by $\varphi_Q$, relabels the momentum $\alpha\mapsto\iota\alpha$ (so $e^{i\alpha\cdot X}$ goes to $e^{i\iota\alpha\cdot X}$) and applies the gauge phase $c(\alpha)$ that compensates the change of Klein factors. For $K=L$ it is the lift of a lattice automorphism to an automorphism of the VOA, i.e.\ a target-space symmetry of the torus theory. For $K\subsetneq L$ it is the inclusion of the sublattice chiral algebra into the larger one.
\end{physics}

\begin{theorem}[Functoriality]\label{ch03:thm:latMorph}
Let $K,L$ be even positive-definite lattices and $\iota\colon K\to L$ an isometry. Then
\begin{enumerate}[label=(\roman*)]
\item $\Phi_\iota$ is a homomorphism of vertex algebras $V_K\to V_L$ (preserves $\mathbb 1$ and all
  products $u_{(n)}v$);
\item if $\operatorname{rank}K=\operatorname{rank}L$, $\Phi_\iota(\omega_K)=\omega_L$ and $\Phi_\iota$
  is injective, with image the set of $x\in V_L$ all of whose momenta lie in $\iota(K)$;
\item $\Phi_\iota\theta_K=\theta_L\Phi_\iota$;
\item $\Phi_\iota g_{\chi\circ\iota}=g_\chi\Phi_\iota$ for every $\chi\colon L\to\Z/2$.
\end{enumerate}
\end{theorem}
\begin{leancode}
theorem latMorph_isVOAHom :
    IsVOAHom (latticeVOAGen K hKe hK) (latticeVOAGen L hLe hL) (latMorph K L ι)
theorem latMorph_mode (u : LatModuleGen K) (n : ℤ) (v : LatModuleGen K) :
    latMorph K L ι (mode (YfullGen K hKe hK u) n v)
      = mode (YfullGen L hLe hL (latMorph K L ι u)) n (latMorph K L ι v)
theorem latMorph_omega : latMorph K L ι (omegaLatGen K) = omegaLatGen L
theorem latMorph_injective : Function.Injective (latMorph K L ι)
theorem mem_range_latMorph_iff (x : LatModuleGen L) :
    x ∈ LinearMap.range (latMorph K L ι) ↔ ∀ γ, γ ∉ Set.range ι → coordEα L γ x = 0
theorem latMorph_thetaLat (x : LatModuleGen K) :
    latMorph K L ι (thetaLat K x) = thetaLat L (latMorph K L ι x)
theorem latMorph_chiLat (χ : L.carrier →+ ZMod 2) (x : LatModuleGen K) :
    latMorph K L ι (chiLat K (χ.comp ι.toAddMonoidHom) x) = chiLat L χ (latMorph K L ι x)
\end{leancode}
(\leanpath{VOA/Moonshine/LatticeMorph.lean}.)

\begin{dictionary}
The hypotheses are section variables, attached by \lean{include} lines in the source:
\begin{itemize}
\item \lean{latMorph_isVOAHom}, \lean{latMorph_mode}: \lean{hK : K.PosDef}, \lean{hL : L.PosDef},
  \lean{hι} (isometry), \lean{hKe : K.Even}, \lean{hLe : L.Even}. \emph{No equal-rank hypothesis.}
\item \lean{latMorph_omega}, \lean{latMorph_injective}, \lean{mem_range_latMorph_iff}:
  \lean{hK}, \lean{hL}, \lean{hι} and \lean{hr : rank K = rank L}.
\item \lean{latMorph_thetaLat}, \lean{latMorph_chiLat}: no hypotheses at all (they are identities
  of linear maps on basis states; \lean{latCob_neg} is unconditional).
\end{itemize}
\lean{χ.comp ι.toAddMonoidHom} is $\chi\circ\iota$. \lean{Set.range ι} is $\iota(K)$, and
\lean{LinearMap.range} the image of a linear map. Since \lean{IsVOAHom} does not include $\omega$
(Section~\ref{ch03:sec:conventions}), statement (i) together with \lean{latMorph_omega} is the full
``conformal homomorphism'' statement in the equal-rank case. The theorem could have been vacuous
if \lean{latMorph} were zero; it is not, since it sends $\mathbb 1$ to $\mathbb 1$ and is injective
in equal rank.
\end{dictionary}

\begin{physics}
This theorem says that lattice isometries act functorially on lattice VOAs: a symmetry of the torus lattice becomes a symmetry of the chiral algebra, compatible with the conformal vector (so it commutes with $L_n$), with the lattice involution $\theta$ ($X\mapsto-X$) and with the sign automorphisms. In particular $\mathrm{Aut}(L)$, the target-space symmetry group of the torus (for the Leech lattice, $\mathrm{Co}_0$), embeds in $\mathrm{Aut}(V_L)$ up to the central extension encoded by the cocycle comparison; a sublattice of equal rank gives the subalgebra of states with momenta in $\iota(K)$, which is the lattice algebra of the smaller (finite-index) lattice. This is what one needs to lift frame automorphisms and sign choices in the Monster construction, but the statement is only the algebraic fact.
\end{physics}

\begin{remark}[Unequal ranks]\label{ch03:rem:notconformal}
If $\operatorname{rank}K<\operatorname{rank}L$, $\Phi_\iota$ is still a vertex algebra
homomorphism, but $\Phi_\iota(\omega_K)$ is the Sugawara vector of the sub-Heisenberg algebra
$\iota(K)\otimes\C$, not $\omega_L$; $\omega_L-\Phi_\iota(\omega_K)$ is the conformal vector of the
orthogonal complement. Injectivity also holds mathematically in that case, but the library proves
it only under \lean{hr} (its proof inverts $\varphi_Q$ by $\varphi_{Q^T}$, which needs
$QQ^T=1$).
\end{remark}

\begin{proof}[Proof on paper]
(i) As for $g_\chi$, check intertwining on generators. Because $\varphi_Q$ is an algebra map with
$Q\alpha^\perp=(\iota\alpha)^\perp$, it sends the Schur polynomials $S_N(\alpha)$ to
$S_N(\iota\alpha)$, and (using $Q^TQ=1$) intertwines $\alpha(n)$ with $(\iota\alpha)(n)$, hence
$E^\pm(\alpha,z)$ with $E^\pm(\iota\alpha,z)$. Then, on $p\otimes e^\beta$,
\begin{equation*}
 \Phi\,Y(e^\alpha,z)(p\otimes e^\beta)=\varepsilon_K(\alpha,\beta)c(\alpha+\beta)z^{\langle\alpha,\beta
 \rangle}\varphi(E^-E^+p)\otimes e^{\iota(\alpha+\beta)}
 =c(\alpha)\,Y(e^{\iota\alpha},z)\,\Phi(p\otimes e^\beta)
\end{equation*}
by Theorem~\ref{ch03:thm:cocCompare} and $\langle\iota\alpha,\iota\beta\rangle=\langle\alpha,
\beta\rangle$. So $\Phi Y(e^\alpha,z)=Y(\Phi e^\alpha,z)\Phi$, as $\Phi e^\alpha=c(\alpha)
e^{\iota\alpha}$. For the currents, $\Phi h_a(z)=(\sum_bQ_{ba}h_b(z))\Phi$, including the zero mode
(this is $Q^T(\iota\beta)^\perp=\beta^\perp$). On the $L$ side, $Y(h_b(-i-1)v,z)={:}\partial^{(i)}
h_b(z)Y(v,z){:}$ for every $v$, and normal ordering is linear in the first factor. Induction over
words gives $\Phi Y(w,z)=Y(\Phi w,z)\Phi$ on word states, hence everywhere.
(ii) $\varphi_Q\omega=\omega$ needs $QQ^T=1$. Injectivity: the momentum components of
$\Phi x$ are $c(\beta)\varphi(x_\beta)$ at $\iota\beta$ and $0$ off $\iota(K)$; $c\neq0$ and
$\varphi$ is invertible. Conversely a vector supported on $\iota(K)$ is the image of
$\sum_\gamma c(\iota^{-1}\gamma)^{-1}\varphi_{Q^T}(x_\gamma)\otimes e^{\iota^{-1}\gamma}$.
(iii) is the computation of Remark~\ref{ch03:rem:Z4} with $c(-\alpha)=c(\alpha)$, plus
$\varphi_Q\sigma=\sigma\varphi_Q$ ($\varphi_Q$ is linear in the variables). (iv) is immediate.
\end{proof}

\begin{status}
All proved; \texttt{\#print axioms} of \lean{latMorph_isVOAHom}, \lean{latMorph_injective},
\lean{mem_range_latMorph_iff}, \lean{latMorph_thetaLat}, \lean{latticeCocycle_map_mul_latCob}
$=$ \texttt{[propext, Classical.choice, Quot.sound]} (re-checked). Conditional only on the
hypotheses listed in the fidelity box. Every statement is in fact proved for the general
\lean{latMorphC ι c} under the minimal hypotheses on $c$ ($c(0)=1$, cocycle relation, $c\ne0$,
$c(-\alpha)=c(\alpha)$ for (iii)) and then specialised.
\end{status}

\subsection{Character automorphisms}

Taking $K=L$, $\iota=\mathrm{id}$ and $c=d$ a character gives automorphisms.

\begin{definition}\label{ch03:def:charAut}
For a homomorphism $d\colon L\to\C^\times$, $g_d(p\otimes e^\alpha)=d(\alpha)\,p\otimes e^\alpha$.
\end{definition}
\noindent File \leanpath{VOA/Moonshine/LatticeMorph.lean}.
\begin{leancode}
noncomputable def charVal (α : L.carrier) : ℂ := ((Additive.toMul (d α) : ℂˣ) : ℂ)

noncomputable def charLat : LatModuleGen L →ₗ[ℂ] LatModuleGen L :=
  latMorphC L L LinearMap.id (charVal L d)

theorem charLat_latStateGen (hL : L.PosDef) (p : FockN (rank L)) (α : L.carrier) :
    charLat L d (latStateGen L p α) = charVal L d α • latStateGen L p α

noncomputable def charAut : Aut (latticeVOAGen L hEven hPos) where ...

theorem charLat_thetaLat_conj (e : L.carrier →+ Additive ℂˣ) (x : LatModuleGen L) :
    charLat L e (thetaLat L (charLat L (-e) x)) = charLat L (2 • e) (thetaLat L x)

theorem exists_half_char : ∃ e : L.carrier →+ Additive ℂˣ, 2 • e = d
\end{leancode}
\begin{dictionary}
\lean{L.carrier →+ Additive ℂˣ} is a homomorphism from $L$ to the multiplicative group $\C^\times$
written additively (Lean's device for viewing a multiplicative group as an additive one); so
\lean{2 • e} is $e^2$ and \lean{-e} is $e^{-1}$, and \lean{charVal} extracts the complex number
$d(\alpha)$. For $d=\mathrm{id}$, $\lean{oscMatrix L L id}=1$ (\lean{oscMatrix_id}), so
\lean{charLat} indeed does not touch the Fock factor. \lean{charLat_thetaLat_conj} is
$g_e\theta g_e^{-1}=g_{e^2}\theta$, needing only \lean{hPos} (to know $Q=1$).
\end{dictionary}

\begin{physics}
A character $d\colon L\to\C^\times$ of the momentum lattice is an element of the dual torus, and $g_d$ multiplies a state of momentum $\alpha$ by $d(\alpha)=e^{i\alpha\cdot\theta}$: it is a translation by $\theta$ on the target torus, i.e.\ an element of the $U(1)^r$ symmetry group generated by the zero modes of the currents. These are the continuous symmetries of the torus CFT present because the weight-one space contains the Cartan currents, and the sign automorphism $g_\chi$ is the special case $d=(-1)^\chi$.
\end{physics}
\begin{proof}[Proof on paper]
$g_d$ is $\Phi^{d}_{\mathrm{id}}$, and $d$ satisfies the cocycle relation with $\varepsilon_L=
\varepsilon_K$ because $d$ is multiplicative. The conjugation formula: on $p\otimes e^\alpha$,
$g_e\theta g_e^{-1}$ multiplies by $e(\alpha)^{-1}e(-\alpha)=e(\alpha)^{-2}$, and $g_{e^2}\theta$
by $e(-\alpha)^2=e(\alpha)^{-2}$. Square roots: $L$ is free, so define $e$ on a basis by choosing
square roots of $d(B_i)$ in $\C$ (algebraically closed), and extend.
\end{proof}

\subsection{Rigidity of lifts of \texorpdfstring{$-1$}{-1}}\label{ch03:sec:rigidity}

FLM show that lifts of $-1\in O(L)$ to $\Aut(V_L)$ are unique up to the characters of $L$. The
library proves the VOA-level statement needed in the Moonshine construction to identify an
abstractly given involution with the standard $\theta$ (the ``$c_\lambda$ residual'' of
\leanpath{Notes/M8GluingPlan.md}).
\noindent File \leanpath{VOA/Moonshine/LatticeMorphRigid.lean}.
\begin{leancode}
noncomputable abbrev heisSt (a : Fin (rank L)) : LatModuleGen L := wordStateLatGen L [(0, a)] 0

noncomputable def hSt (α : L.carrier) : LatModuleGen L := latStateGen L (vSchur L α 1) 0
\end{leancode}
\lean{heisSt L a} $=h_a(-1)\mathbb 1$ (frame direction), and \lean{hSt L α} $=\alpha(-1)\mathbb 1$
since $S_1(\alpha)=v_0(\alpha)=\sum_a\alpha^\perp_aX_{(0,a)}$ (\lean{hSt} lives in
\leanpath{VOA/Moonshine/TrialitySl2.lean}).

\begin{theorem}[Rigidity]\label{ch03:thm:rigidity}
Let $L$ be even positive-definite and $\psi\in\Aut(V_L)$ such that $\psi(\alpha(-1)\mathbb 1)=
-\alpha(-1)\mathbb 1$ for all $\alpha\in L$ (equivalently $\psi(h_a(-1)\mathbb 1)=-h_a(-1)\mathbb1$
for all $a$) and $\psi(\mathcal F\otimes e^\gamma)\subseteq\mathcal F\otimes e^{-\gamma}$ for all
$\gamma$. Then $\psi=g_d\theta$ for a character $d\colon L\to\C^\times$, and $\psi=g_e\theta
g_e^{-1}$ for a character $e$; in particular $\psi$ is conjugate to $\theta$ in $\Aut(V_L)$.
\end{theorem}
\begin{leancode}
theorem exists_char_of_lift (ψ : Aut (latticeVOAGen L hEven hPos))
    (hh : ∀ a, ψ (heisSt L a) = -heisSt L a)
    (hsec : ∀ (γ : L.carrier) (x : LatModuleGen L), x ∈ chargeSec L γ → ψ x ∈ chargeSec L (-γ)) :
    ∃ d : L.carrier →+ Additive ℂˣ, ∀ x, ψ x = charLat L d (thetaLat L x)

theorem lift_eq_conj_theta (ψ : Aut (latticeVOAGen L hEven hPos))
    (hh : ∀ a, ψ (heisSt L a) = -heisSt L a)
    (hsec : ∀ (γ : L.carrier) (x : LatModuleGen L), x ∈ chargeSec L γ → ψ x ∈ chargeSec L (-γ)) :
    ∃ e : L.carrier →+ Additive ℂˣ, ∀ x, ψ x = charLat L e (thetaLat L (charLat L (-e) x))

theorem lift_eq_conj_theta' (ψ : Aut (latticeVOAGen L hEven hPos))
    (hh : ∀ α, ψ (hSt L α) = -hSt L α)
    (hsec : ∀ (γ : L.carrier) (x : LatModuleGen L), x ∈ chargeSec L γ → ψ x ∈ chargeSec L (-γ)) :
    ∃ e : L.carrier →+ Additive ℂˣ, ∀ x, ψ x = charLat L e (thetaLat L (charLat L (-e) x))
\end{leancode}
(\leanpath{VOA/Moonshine/LatticeMorphRigid.lean}.)

\begin{dictionary}
\lean{chargeSec L γ} is $\mathcal F\otimes e^\gamma$. The two hypotheses say that $\psi$ acts as
$-1$ on $\mathfrak h=(V_L)_1^{\text{Heis}}$ and maps momentum $\gamma$ to momentum $-\gamma$; this is
what ``$\psi$ lifts $-1\in O(L)$'' means. The primed version states the first hypothesis on the
lattice-meaningful states $\alpha(-1)\mathbb 1$, which span the same space as the $h_a(-1)\mathbb1$:
$\alpha(-1)\mathbb1=\sum_a\alpha^\perp_ah_a(-1)\mathbb1$, and the vectors $B_i^\perp=CB_i$ span
$\chthreeR^r$ because $C=\sqrt G$ is invertible (the library proves the needed inclusion as
\lean{heisSt_mem_span_hSt}); it is the form used downstream.
The conclusion is an equality of functions, quantified over all $x$. The sector hypothesis is an honest extra input: mathematically it follows from the first one
(the momentum sectors are the joint eigenspaces of the zero modes $h_a(0)$, and $\psi$ conjugates
$h(0)$ to $-h(0)$), but the library does not prove that implication, so the theorem assumes it.
\end{dictionary}

\begin{physics}
This is a statement that every automorphism acting as $-1$ on the oscillator currents and swapping each momentum sector $\gamma\to-\gamma$ is conjugate to the reflection $\theta$ ($X\mapsto-X$) by a translation in the torus (a $U(1)^r$ symmetry). Physically, once the current algebra is fixed to be acted on by $-1$, the only freedom left is a shift of the origin of the torus, so such reflection symmetries (the $\Z/2$ target-space inversion with fixed points) form a single conjugacy class. This is the rigidity input used to identify the lifted involution with the standard one in the orbifold construction.
\end{physics}

\begin{proof}[Proof on paper]
Let $\phi=\psi\theta$. It fixes every $h_a(-1)\mathbb 1$ (two sign changes), preserves each sector
$\mathcal F\otimes e^\gamma$, and, being an automorphism, preserves the $L_0$-grading. Hence
$\phi(e^\gamma)$ lies in $(\mathcal F\otimes e^\gamma)\cap(V_L)_{|\gamma|^2/2}=\mathcal F_0\otimes
e^\gamma=\C e^\gamma$; write $\phi(e^\gamma)=d(\gamma)e^\gamma$. Since $(h_a(-1)\mathbb
1)_{(-i-1)}=h_a(-i-1)$ and $\phi$ fixes $h_a(-1)\mathbb 1$, $\phi$ commutes with all creation
oscillators, so $\phi=d(\gamma)$ on the whole sector. Applying $\phi$ to
$(e^\alpha)_{(-\langle\alpha,\beta\rangle-1)}e^\beta=\varepsilon(\alpha,\beta)e^{\alpha+\beta}\neq0$
gives $d(\alpha)d(\beta)=d(\alpha+\beta)$; $d(\gamma)\neq0$ by injectivity. Thus $\psi=\phi\theta^{-1}
=g_d\theta$. Writing $d=e^2$ (\lean{exists_half_char}) and using $g_e\theta g_e^{-1}=g_{e^2}\theta$
gives the conjugacy.
\end{proof}

\begin{status}
Proved, conditional only on \lean{hEven}, \lean{hPos} and the two displayed hypotheses on $\psi$;
\texttt{\#print axioms exists\_char\_of\_lift}, \lean{lift_eq_conj_theta'} $=$
\texttt{[propext, Classical.choice, Quot.sound]} (re-checked).
\end{status}

Finally, an automorphism of $V_L$ preserving the image of an injective conformal homomorphism
$\Phi\colon V\to W$ pulls back to an automorphism of $V$:
\begin{leancode}
noncomputable def Aut.pullback (hΦ : IsVOAHom V W Φ)
    (hω : Φ V.conformalVector = W.conformalVector) : Aut V
theorem Aut.pullback_apply (hΦ : IsVOAHom V W Φ)
    (hω : Φ V.conformalVector = W.conformalVector) (x : 𝒱) :
    Φ (Aut.pullback Φ hinj g hg hΦ hω x) = g (Φ x)
\end{leancode}
with section variables \lean{hinj : Function.Injective Φ}, \lean{g : Aut W} and
\lean{hg : (LinearMap.range Φ).map g.toLinearEquiv.toLinearMap = LinearMap.range Φ}. Note the
explicit hypothesis \lean{hω}, needed precisely because \lean{IsVOAHom} omits $\omega$.

\subsection{Where functoriality is used}

The first application (\leanpath{VOA/Moonshine/FrameBridge.lean}) takes $K=\Leech_0\cong A$, an
index-$2$ sublattice common to the Leech lattice and the Niemeier lattice $N(A_1^{24})$ (the
even unimodular lattice of rank $24$ whose roots form $24$ orthogonal copies of $A_1$, obtained
from $\sqrt2\,\Z^{24}$ by gluing along the Golay code; Chapters~\ref{ch:02} and~\ref{ch:11}), with the
two isometric embeddings $\iota_\Leech\colon a\mapsto 2a$ and $\iota_N\colon A\subseteq N$. The
resulting \lean{bridgeΛ := latMorph lam0 leech inclΛ} and \lean{bridgeN := latMorph lam0 golayN inclN}
are injective conformal homomorphisms commuting with $\theta$, and give the vertex algebra
isomorphism $V_{\Leech_0}^+\cong V_0^+$ (\lean{lam0PlusEquiv}) through which the triality
automorphism of $V_{N}$ (Chapter~\ref{ch:11}) is transported to $V_\Leech$. This is why both the
exact $\theta$-compatibility (Remark~\ref{ch03:rem:Z4}) and the rigidity theorem are needed.

\section{The design question: a non-computable basis}\label{ch03:sec:design}

The cocycle \lean{latticeCocycle} is defined over \lean{Module.Free.chooseBasis}, a basis selected
by the axiom of choice, and the oscillator frame is $\sqrt G$ in that basis
(Sections~\ref{ch03:sec:cocycle}--\ref{ch03:sec:frame}). Consequently \emph{no explicit sign}
$\varepsilon(\alpha,\beta)$ and no explicit frame coordinate can be computed in Lean: there is no
way to evaluate, say, $\varepsilon_\Leech(\alpha,\beta)$ for two given Leech vectors, nor to say
which Fock colour corresponds to which Golay coordinate.

During the construction of the triality automorphism (the step called R-top in
\leanpath{Notes/M8GluingPlan.md}, Status 2026-10-07 ``R-top partial''), this was recorded as a
\emph{design issue}: ``the Leech cocycle is \lean{latticeCocycle} over
\lean{Module.Free.chooseBasis} and the oscillator frame is $\sqrt G$, not the Golay frame --- NO
explicit sign is computable; every identity must stay symbolic. An explicit (Golay-frame)
basis/cocycle model, transported via the functoriality machinery (\lean{latMorph}), could
drastically simplify R-top's remaining work; decision pending.'' The question was put to both
maintainers.

The question was resolved without a refactor. The later status entry (``$\sigma$ EXISTS on both
twisted sectors --- unconditional'') records: \emph{``DESIGN QUESTION RESOLVED: \ldots an
explicit-basis refactor would NOT have materially reduced cost (all identities reduced to
bimultiplicative-cocycle algebra / $\xi$ / Golay facts; no concrete cocycle value on
\lean{chooseBasis} was ever needed). Keep the foundations as they are.''} The reasons, in
mathematical terms:
\begin{enumerate}
\item Every identity that the Moonshine construction needs about $\varepsilon$ follows from
  bimultiplicativity, $\varepsilon^2=1$ and the commutator property
  $\varepsilon(\alpha,\beta)\varepsilon(\beta,\alpha)=(-1)^{\langle\alpha,\beta\rangle}$. These
  determine $\varepsilon$ up to coboundary, and the vertex operator algebra up to isomorphism; no
  specific value is ever required.
\item Where two different cocycles had to be compared (two lattices, two bases) the functoriality
  machinery of Section~\ref{ch03:sec:functoriality} does it once and for all, with the exactly
  $\theta$-compatible $\Z/4$ coboundary.
\item Where an unknown scalar appeared (e.g.\ the $c_\lambda$ in $\Theta e^\lambda=c_\lambda
  e^{-\lambda}$), rigidity (Theorem~\ref{ch03:thm:rigidity}) shows it is a character and can be
  conjugated away; one never needs its value.
\item The frame directions only enter through the orthogonality relations $Q^TQ=QQ^T=1$ and the
  identity $Q\alpha^\perp=(\iota\alpha)^\perp$, never through explicit coordinates.
\end{enumerate}

\begin{warning}
For a referee the relevant lesson is the converse: a statement in this part of the library that
\emph{appears} to depend on a specific value of \lean{latticeCocycle} (or of \lean{orthCoord})
should be read with suspicion, because such a value cannot be computed and any proof of it must
come from basis-independent properties. No such statement is used in the construction of
$\Vnat$ (Chapter~\ref{ch:14}).
\end{warning}

\section{Referee's checklist}\label{ch03:sec:checklist}

Abbreviations: E $=$ \lean{hEven} (evenness), P $=$ \lean{hPos} (positive-definiteness), I $=$
isometry \lean{hι}, R $=$ equal rank \lean{hr}. ``Uncond.'' means no hypotheses beyond the
displayed data; all entries have \texttt{\#print axioms} $\subseteq$ \texttt{[propext,
Classical.choice, Quot.sound]} and none uses \lean{VOA.pbw}. Namespaces are omitted in the table: \lean{VOA.HeisenbergN} for the first three rows, \lean{VOA.IntegralLattice} for the cocycle and frame rows, \lean{VOA.Reconstruction} for \lean{reconstruct}, \lean{VOA.Moonshine} for the \lean{leech...} rows, and \lean{VOA.LatGen} otherwise.

{\footnotesize
\begin{longtable}{>{\raggedright\arraybackslash}p{5.0cm}>{\raggedright\arraybackslash}p{4.7cm}>{\raggedright\arraybackslash}p{2.6cm}>{\raggedright\arraybackslash}p{1.6cm}}
\toprule
Lean declaration & Intended meaning & Hypotheses & Fidelity\\
\midrule\endhead
\lean{oscN_comm} & $[h_a(m),h_b(n)]=mk\delta_{ab}\delta_{m+n,0}$ on the Fock polynomial ring & uncond. & Thm.~\ref{ch03:thm:oscN_comm}\\
\lean{L_virasoroN} & Sugawara modes satisfy Virasoro with $c=r$ & $k\ne0$ & Thm.~\ref{ch03:thm:virasoroN}\\
\lean{heisenbergNVOA} & rank-$r$ free boson VOA, CFT type, $c=r$ & $k\neq0$ & Thm.~\ref{ch03:thm:heisVOA}\\
\lean{cocycle_cocycleLaw} & $\varepsilon_{\mathcal B}$ is a 2-cocycle (any basis) & uncond. & Thm.~\ref{ch03:thm:cocycle}\\
\lean{latticeCocycle_comm} & $\varepsilon(\alpha,\beta)\varepsilon(\beta,\alpha)=(-1)^{\langle\alpha,\beta\rangle}$ & E & Thm.~\ref{ch03:thm:cocycle}\\
\lean{orthCoord_inner} & $\alpha^\perp\cdot\beta^\perp=\langle\alpha,\beta\rangle$ & P & Def.~\ref{ch03:def:frame}\\
\lean{chargeOsc_comm} & $[\alpha(m),\beta(n)]=\langle\alpha,\beta\rangle m\delta_{m+n,0}$ & P & Def.~\ref{ch03:def:chargeOsc}\\
\lean{gradingLatGen_isInternal_of_posDef}, \lean{gradingLatGen_finrank_of_posDef} & $V_L=\bigoplus(V_L)_n$, finite-dim., CFT type & E, P & Thm.~\ref{ch03:thm:gradingLat}\\
\lean{Yeαvec_tmul} & defining formula of $Y(e^\alpha,z)$ with cocycle & uncond. & Def.~\ref{ch03:def:Yealpha}\\
\lean{Yeαvec_local} & $Y(e^\alpha),Y(e^\beta)$ mutually local & E, P & Thm.~\ref{ch03:thm:Yelocal}\\
\lean{reconstruct} & existence of $Y$ from local covariant generators & as listed & Thm.~\ref{ch03:thm:reconstruct}\\
\lean{latYfullGen_of_locality} & the vertex map of $V_L$ (unique) & E, P & Thm.~\ref{ch03:thm:latYfull}\\
\lean{latticeVOAGen} & $V_L$ is a VOA of CFT type & E, P & Thm.~\ref{ch03:thm:latticeVOA}\\
\lean{latticeVOAGen_centralCharge} & $c(V_L)=\operatorname{rank}L$ & E, P & Thm.~\ref{ch03:thm:latticeVOA}\\
\lean{leechVOA} & $V_\Leech$ & uncond. & Def.~\ref{ch03:def:leechVOA}\\
\lean{leechVOA_centralCharge} & $c(V_\Leech)=24$ & uncond. & Def.~\ref{ch03:def:leechVOA}\\
\lean{gradingLatGen_one_of_rootless} & rootless $\Rightarrow (V_L)_1=\mathfrak h(-1)\mathbb 1$ & E, P, no roots & Thm.~\ref{ch03:thm:rootless}\\
\lean{leechVOA_grading_one} & $(V_\Leech)_1=\mathfrak h(-1)\mathbb1$ & uncond. & Thm.~\ref{ch03:thm:rootless}\\
\lean{chiAut}, \lean{chiAut_sq} & $g_\chi\in\Aut V_L$, $g_\chi^2=1$ & E, P & Thm.~\ref{ch03:thm:chiAut}\\
\lean{chiAut_mul_thetaAut} & $g_\chi\theta=\theta g_\chi$ & E, P & Thm.~\ref{ch03:thm:chiAut}\\
\lean{chiLat_eq_self_iff} & $g_\chi$-fixed $\Leftrightarrow$ momenta in $\ker\chi$ & uncond. & Thm.~\ref{ch03:thm:chieig}\\
\lean{cobSign_add} & $\Z/4$ coboundary identity & $D+D^T$ even & Lem.~\ref{ch03:lem:cobSign}\\
\lean{latticeCocycle_map_mul_latCob} & $\varepsilon_L\circ\iota$ and $\varepsilon_K$ cohomologous via even $c$ & E (both), I & Thm.~\ref{ch03:thm:cocCompare}\\
\lean{latMorph_isVOAHom} & $\Phi_\iota$ is a VA homomorphism & E, P (both), I & Thm.~\ref{ch03:thm:latMorph}\\
\lean{latMorph_omega} & $\Phi_\iota\omega=\omega$ & P, I, R & Thm.~\ref{ch03:thm:latMorph}\\
\lean{latMorph_injective} & $\Phi_\iota$ injective & P, I, R & Thm.~\ref{ch03:thm:latMorph}\\
\lean{mem_range_latMorph_iff} & image $=$ momenta in $\iota(K)$ & P, I, R & Thm.~\ref{ch03:thm:latMorph}\\
\lean{latMorph_thetaLat} & $\Phi_\iota\theta=\theta\Phi_\iota$ exactly & uncond. & Thm.~\ref{ch03:thm:latMorph}, Rem.~\ref{ch03:rem:Z4}\\
\lean{charAut} & $g_d\in\Aut V_L$ for $d\colon L\to\C^\times$ & E, P & Def.~\ref{ch03:def:charAut}\\
\lean{exists_char_of_lift} & lift of $-1$ is $g_d\theta$ & E, P, two hyps on $\psi$ & Thm.~\ref{ch03:thm:rigidity}\\
\lean{lift_eq_conj_theta'} & lift of $-1$ is conjugate to $\theta$ & E, P, two hyps on $\psi$ & Thm.~\ref{ch03:thm:rigidity}\\
\bottomrule
\end{longtable}}

Points a referee should specifically confirm:
\begin{enumerate}
\item \emph{Indexing.} \lean{FormalDist} is indexed by powers of $z$; modes via
  \lean{mode f n = f (-n-1)}. The Fock variable \lean{X (i,a)} is $h_a(-i-1)\mathbb1$.
\item \emph{Naming.} The library's $E^\mp(\alpha,z)$ are FLM's $E^\mp(-\alpha,z)$.
\item \emph{Choice.} \lean{YfullGen} is defined by \lean{.choose} but is unique; the cocycle is
  defined via \lean{chooseBasis} and is non-computable but satisfies all needed identities.
\item \emph{Homomorphisms.} \lean{IsVOAHom} omits $\omega$; conformality is
  \lean{latMorph_omega}, under equal rank.
\item \emph{The $\Z/4$ coboundary} is required for exact $\theta$-compatibility
  (Remark~\ref{ch03:rem:Z4}); the $\pm1$ description in one module docstring is stale.
\item \emph{Stale docstrings} (no effect on statements): \leanpath{LatticeGenGrading.lean} calls a
  discharged lemma a crux; \leanpath{LatticeGenModule.lean} lists ``remaining Goals'' that are
  done; \leanpath{LatticeMorphFock.lean} mentions a nonexistent \lean{exists_fockMatrix}.
\end{enumerate}

\section*{Notes and references}\addcontentsline{toc}{section}{Notes and references}
\label{ch03:sec:notes}

The free boson and its Fock space are classical; the vertex-algebraic treatment follows
I.~B.~Frenkel, J.~Lepowsky and A.~Meurman, \emph{Vertex Operator Algebras and the Monster}, Pure
and Applied Mathematics~134, Academic Press, 1988, Chapters~1--4 (Heisenberg algebras and their
Fock spaces) and Chapters~5--8 (central extensions of lattices, the cocycle
$\varepsilon$, the vertex operators $Y(e^\alpha,z)=E^-(-\alpha,z)E^+(-\alpha,z)e_\alpha z^\alpha$,
and the lattice vertex operator algebra $V_L$); our Definition~\ref{ch03:def:Yealpha} is their
(8.4.1) up to the naming of $E^\pm$. The lifts of $-1$ and their uniqueness up to characters are in
FLM Chapter~10 and Proposition~5.4.1 there.

The vertex operator realisation of affine Lie algebras by lattice vertex operators is due to
I.~B.~Frenkel and V.~G.~Kac, ``Basic representations of affine Lie algebras and dual resonance
models'', \emph{Invent.\ Math.}~62 (1980), 23--66, and G.~Segal, ``Unitary representations of some
infinite dimensional groups'', \emph{Comm.\ Math.\ Phys.}~80 (1981), 301--342. The first general
definition of a vertex algebra, with lattice vertex algebras as the main example, is
R.~E.~Borcherds, ``Vertex algebras, Kac--Moody algebras, and the Monster'', \emph{Proc.\ Nat.\
Acad.\ Sci.\ USA}~83 (1986), 3068--3071.

The axiomatics by locality and the reconstruction (``existence'') theorem used in
Section~\ref{ch03:sec:reconstruction} are from V.~G.~Kac, \emph{Vertex Algebras for Beginners},
2nd ed., University Lecture Series~10, AMS, 1998, \S\S4.4--4.5 and \S5.4--5.5 (lattice vertex
algebras); see also E.~Frenkel and D.~Ben-Zvi, \emph{Vertex Algebras and Algebraic Curves},
2nd ed., Mathematical Surveys and Monographs~88, AMS, 2004, Chapters~2 and~5 (free bosons,
lattice vertex algebras, reconstruction), and J.~Lepowsky and H.~Li, \emph{Introduction to Vertex
Operator Algebras and Their Representations}, Progress in Mathematics~227, Birkh\"auser, 2004,
Chapter~6 (Heisenberg and lattice vertex operator algebras via the Jacobi identity), where
the cocycle normalisation of Section~\ref{ch03:sec:cocycle} is also discussed. Dong's lemma is from
C.~Dong, ``Vertex algebras associated with even lattices'', \emph{J.\ Algebra}~161 (1993),
245--265, and is presented in Kac's book (Lemma~3.2). Goddard's uniqueness theorem is from
P.~Goddard, ``Meromorphic conformal field theory'', in \emph{Infinite-dimensional Lie algebras and
groups} (Luminy--Marseille, 1988), Adv.\ Ser.\ Math.\ Phys.~7, World Scientific, 1989, 556--587.
The library's own primary reference for the definitions is T.~Gannon, \emph{Moonshine beyond the
Monster}, Cambridge University Press, 2006, \S5.2.

The Leech lattice and its construction from the Golay code are in W.~Ebeling, \emph{Lattices and
Codes}, 3rd ed., Springer Spektrum, 2013, and J.~H.~Conway and N.~J.~A.~Sloane, \emph{Sphere
Packings, Lattices and Groups}, 3rd ed., Springer, 1999; uniqueness is J.~H.~Conway, ``A
characterisation of Leech's lattice'', \emph{Invent.\ Math.}~7 (1969), 137--142 (not formalised).
The use of the sublattice $\Leech_0$ and of the Niemeier lattice $N(A_1^{24})$ to construct the
triality follows C.~Dong, R.~L.~Griess Jr.\ and G.~H\"ohn, ``Framed vertex operator algebras,
codes and the Moonshine module'', \emph{Comm.\ Math.\ Phys.}~193 (1998), 407--448
(arXiv:q-alg/9707008); see Chapters~\ref{ch:11}--\ref{ch:13}.

The design discussion of Section~\ref{ch03:sec:design} and the $\Z/4$ coboundary are recorded in
the repository files \leanpath{Notes/M8GluingPlan.md} (Status sections of 2026-10-06 and
2026-10-07).
\endgroup
